\documentclass[12pt, letterpaper]{article}
\usepackage[utf8]{inputenc}
\usepackage{graphicx}
\usepackage{fullpage}
\usepackage{setspace}
\usepackage{biblatex}
\usepackage{mathtools}
\usepackage{appendix}
\usepackage[version=4]{mhchem}
\usepackage[export]{adjustbox}
\usepackage{pdflscape}
\usepackage{listings}
\usepackage[hidelinks]{hyperref}
\usepackage{pdfpages}
\addbibresource{references.bib}
\doublespacing{}
\begin{document}


\section{Residual Helmholz Energy}

% EqID: ares_total
% Source: No explicit citation in equations.tex context
% Status: Project summary relation
% Description: Provides a residual Helmholtz-energy relation for residual helmholz energy.
% Change note: No explicit citation on this equation block in the source file.

\begin{align}
    \tilde{a}^{\mathrm{res}}=\tilde{a}^{h c}+\tilde{a}^{\text {disp }}+\tilde{a}^{\text {assoc }}+\tilde{a}^{\text {DH }}+\tilde{a}^{\text {Born }}
    \label{eq:ares_total}
\end{align}

\subsection{Hard-Chain Reference Contribution}


See~\cite{Gross2001}.
% EqID: ares_hc
% Source: \cite{Gross2001}, Eq.~(A.4)
% Status: Close literature match
% Description: Provides a residual Helmholtz-energy relation for hard-chain reference contribution.
% Change note: High textual similarity to a tagged equation in the cited local paper export.

\begin{equation}
    \tilde{a}^{\mathrm{hc}}=\bar{m} \tilde{a}^{\mathrm{hs}}-\sum_{i} x_{i}\left(m_{i}-1\right) \ln \mathrm{g}_{i i}^{\mathrm{hs}}\left(\sigma_{i i}\right)
    \label{eq:ares_hc}
\end{equation}

% EqID: m_bar
% Source: \cite{Gross2001}, Eq.~(A.5)
% Status: Close literature match
% Description: Provides a supporting relation used in hard-chain reference contribution.
% Change note: High textual similarity to a tagged equation in the cited local paper export.

\begin{equation}
    \bar{m}=\sum_{i} x_{i} m_{i}
    \label{eq:m_bar}
\end{equation}

% EqID: ares_hs
% Source: \cite{Gross2001}, Eq.~(A.6)
% Status: Close literature match
% Description: Provides a residual Helmholtz-energy relation for hard-chain reference contribution.
% Change note: High textual similarity to a tagged equation in the cited local paper export.

\begin{equation}
    \tilde{a}^{\mathrm{hs}} =\frac{1}{\zeta_{0}}\left[\frac{3 \zeta_{1} \zeta_{2}}{\left(1-\zeta_{3}\right)}+\frac{\zeta_{2}^3}{\zeta_{3}\left(1-\zeta_{3}\right)^2}+\left(\frac{\zeta_{2}^3}{\zeta_{3}^2}-\zeta_{0}\right) \ln \left(1-\zeta_{3}\right)\right]
    \label{eq:ares_hs}
\end{equation}

% EqID: g_hs_contact
% Source: \cite{Gross2001}, Eq.~(A.7)
% Status: Close literature match
% Description: Provides a supporting relation used in hard-chain reference contribution.
% Change note: High textual similarity to a tagged equation in the cited local paper export.

\begin{equation}
    \mathrm{g}_{i j}^{\mathrm{hs}}=\frac{1}{\left(1-\zeta_{3}\right)}+\left(\frac{d_{i} d_{j}}{d_{i}+d_{j}}\right) \frac{3 \zeta_{2}}{\left(1-\zeta_{3}\right)^2} + \left(\frac{d_{i} d_{j}}{d_{i}+d_{j}}\right)^2 \frac{2 \xi_{2}{ }^2}{\left(1-\xi_{3}\right)^3}
    \label{eq:g_hs_contact}
\end{equation}

% EqID: zeta_n
% Source: \cite{Gross2001}, Eq.~(A.8)
% Status: Close literature match
% Description: Provides a supporting relation used in hard-chain reference contribution.
% Change note: High textual similarity to a tagged equation in the cited local paper export.

\begin{equation}
    \xi_{n}=\frac{\pi}{6} \rho \sum_{i} x_{i} m_{i} d_{i}^n \quad n \in\{0,1,2,3\}
    \label{eq:zeta_n}
\end{equation}

% EqID: d_segment
% Source: \cite{Gross2001}, Eq.~(A.9)
% Status: Close literature match
% Description: Provides a supporting relation used in hard-chain reference contribution.
% Change note: High textual similarity to a tagged equation in the cited local paper export.

\begin{equation}
    d_{i}=\sigma_{i}\left[1-0.12 \exp \left(-3 \frac{\epsilon_{i}}{k T}\right)\right]
    \label{eq:d_segment}
\end{equation}

\subsection{Dispersion Contribution}

See~\cite{Gross2001}.
% EqID: ares_disp
% Source: \cite{Gross2001}, Eq.~(A.10)
% Status: Close literature match
% Description: Provides a residual Helmholtz-energy relation for dispersion contribution.
% Change note: High textual similarity to a tagged equation in the cited local paper export.

\begin{equation}
    \tilde{a}^{\mathrm{disp}}=-2 \pi \rho I_{1}(\eta, \bar{m}) \overline{m^2 \epsilon \sigma^3}-\pi \rho \bar{m} C_{1} I_{2}(\eta, \bar{m}) \overline{m^2 \epsilon^2 \sigma^3}
    \label{eq:ares_disp}
\end{equation}

% EqID: c1_disp
% Source: \cite{Gross2001}, Eq.~(A.11)
% Status: Adapted notation
% Description: Provides a residual Helmholtz-energy relation for dispersion contribution.
% Change note: Moderate-to-high similarity; notation/arrangement appears adapted from the cited equation.

\begin{equation}
    C_{1} = \left(1+\bar{m} \frac{8 \eta-2 \eta^2}{(1-\eta)^4}+\right.\left.\quad(1-\bar{m}) \frac{20 \eta-27 \eta^2+12 \eta^3-2 \eta^4}{[(1-\eta)(2-\eta)]^2}\right)
    \label{eq:c1_disp}
\end{equation}

% EqID: m2epssigma3_bar
% Source: \cite{Gross2001}, Eq.~(A.13)
% Status: Adapted implementation form
% Description: Provides a residual Helmholtz-energy relation for dispersion contribution.
% Change note: Lower similarity; likely algebraically adapted for implementation or combined terms.

\begin{equation}
    \overline{m^2 \epsilon \sigma^3}=\sum_{i} \sum_{j} x_{i} x_{j} m_{i} m_{j}\left(\frac{\epsilon_{i j}}{k T}\right) \sigma_{i j}^3
    \label{eq:m2epssigma3_bar}
\end{equation}

% EqID: m2eps2sigma3_bar
% Source: \cite{Gross2001}, Eq.~(A.13)
% Status: Adapted implementation form
% Description: Provides a residual Helmholtz-energy relation for dispersion contribution.
% Change note: Lower similarity; likely algebraically adapted for implementation or combined terms.

\begin{equation}
    \overline{m^2 \epsilon^2 \sigma^3}=\sum_{i} \sum_{j} x_{i} x_{j} m_{i} m_{j}\left(\frac{\epsilon_{i j}}{k T}\right)^2 \sigma_{i j}{ }^3
    \label{eq:m2eps2sigma3_bar}
\end{equation}

% EqID: sigma_ij
% Source: \cite{Gross2001}, Eq.~(A.15)
% Status: Adapted notation
% Description: Provides a supporting relation used in dispersion contribution.
% Change note: Moderate-to-high similarity; notation/arrangement appears adapted from the cited equation.

\begin{equation}
    \sigma_{i j}=\frac{1}{2}\left(\sigma_{i}+\sigma_{j}\right) \cdot \left(1- l_{ij} \right)
    \label{eq:sigma_ij}
\end{equation}

% EqID: epsilon_ij_mixing
% Source: \cite{Gross2001}, Eq.~(A.15)
% Status: Project-specific modification
% Description: Provides a supporting relation used in dispersion contribution.
% Change note: Mapped to Eq. (A.15); this file adds a piecewise ion-specific zero-dispersion branch not written in the original equation.

\begin{equation}
    \epsilon_{ij}=
    \begin{cases}
        \sqrt{\epsilon_{i}\,\epsilon_{j}}\left(1-k_{ij}\right), & z_{i} z_{j} \neq 0, \\[6pt]
        0,                                                  & z_{i} z_{j} = 0 .
    \end{cases}
    \label{eq:epsilon_ij_mixing}
\end{equation}

% EqID: epsilon_ij_ionic_zero
% Source: \cite{Gross2001}, Eq.~(A.15)
% Status: Project-specific modification
% Description: Provides a supporting relation used in dispersion contribution.
% Change note: Represents the zero-dispersion limiting branch of the Eq. (A.15) combining rule for selected interaction classes.

\begin{equation}
    \epsilon_{i j}= 0
    \label{eq:epsilon_ij_ionic_zero}
\end{equation}

% EqID: i1_disp
% Source: \cite{Gross2001}, Eq.~(A.17)
% Status: Adapted implementation form
% Description: Provides a residual Helmholtz-energy relation for dispersion contribution.
% Change note: Lower similarity; likely algebraically adapted for implementation or combined terms.

\begin{equation}
    I_{1}(\eta, \bar{m})=\sum_{i=0}^6 a_{i}(\bar{m}) \eta^i
    \label{eq:i1_disp}
\end{equation}

% EqID: i2_disp
% Source: \cite{Gross2001}, Eq.~(A.17)
% Status: Adapted implementation form
% Description: Provides a residual Helmholtz-energy relation for dispersion contribution.
% Change note: Lower similarity; likely algebraically adapted for implementation or combined terms.

\begin{equation}
    I_{2}(\eta, \bar{m})=\sum_{i=0}^6 b_{i}(\bar{m}) \eta^i
    \label{eq:i2_disp}
\end{equation}

% EqID: a_i_mbar
% Source: \cite{Gross2001}, Eq.~(A.19)
% Status: Adapted implementation form
% Description: Provides a supporting relation used in dispersion contribution.
% Change note: Lower similarity; likely algebraically adapted for implementation or combined terms.

\begin{equation}
    a_{i}(\bar{m})=a_{0 i}+\frac{\bar{m}-1}{\bar{m}} a_{1 i}+\frac{\bar{m}-1}{\bar{m}} \frac{\bar{m}-2}{\bar{m}} a_{2 i}
    \label{eq:a_i_mbar}
\end{equation}

% EqID: b_i_mbar
% Source: \cite{Gross2001}, Eq.~(A.19)
% Status: Adapted implementation form
% Description: Provides a supporting relation used in dispersion contribution.
% Change note: Lower similarity; likely algebraically adapted for implementation or combined terms.

\begin{equation}
    b_{i}(\bar{m})=b_{0 i}+\frac{\bar{m}-1}{\bar{m}} b_{1 i}+\frac{\bar{m}-1}{\bar{m}} \frac{\bar{m}-2}{\bar{m}} b_{2 i}
    \label{eq:b_i_mbar}
\end{equation}

% EqID: rho_from_eta
% Source: \cite{Gross2001}, Eq.~(A.20)
% Status: Close literature match
% Description: Provides a supporting relation used in dispersion contribution.
% Change note: High textual similarity to a tagged equation in the cited local paper export.

\begin{equation}
    \rho=\frac{6}{\pi} \eta\left(\sum_{i} x_{i} m_{i} d_{i}^3\right)^{-1}
    \label{eq:rho_from_eta}
\end{equation}

% EqID: rho_reduced
% Source: \cite{Gross2001}, Eq.~(A.21)
% Status: Close literature match
% Description: Provides a supporting relation used in dispersion contribution.
% Change note: High textual similarity to a tagged equation in the cited local paper export.

\begin{equation}
    \hat{\rho}=\frac{\rho}{N_{\mathrm{AV}}}\left(10^{10} \frac{A^\circ{}}{\mathrm{~m}}\right)^3\left(10^{-3} \frac{\mathrm{kmol}}{\mathrm{~mol}}\right)
    \label{eq:rho_reduced}
\end{equation}

\subsection{Association Contribution}

See~\cite{Chapman1990}.
% EqID: ares_assoc
% Source: \cite{Chapman1990} (exact equation number not available in local PDFs)
% Status: No direct numbered source in local corpus
% Description: Provides a residual Helmholtz-energy relation for association contribution.
% Change note: Association Helmholtz form is traced to Chapman/Wertheim SAFT association theory, but the exact numbered equation is not present in the local progression PDFs.

\begin{equation}
    \tilde{a}^{\mathrm{assoc}}= \sum_{i} x_{i}\sum_{\mathrm{A}_{i}}\left(\ln X^{\mathrm{A}_{i}}-\frac{X^{\mathrm{A}_{i}}}{2}+\frac{1}{2}\right)
    \label{eq:ares_assoc}
\end{equation}

% EqID: x_assoc_site
% Source: \cite{Chapman1990} (exact equation number not available in local PDFs)
% Status: No direct numbered source in local corpus
% Description: Provides a residual Helmholtz-energy relation for association contribution.
% Change note: Mass-action site-fraction relation from SAFT association theory; exact numbered equation unavailable in the local PDF set.

\begin{equation}
    \begin{aligned}
        X^{\mathrm{A}_{i}}
        &=\left[1+ \sum_{j} \sum_{\mathrm{B}_{j}} \rho x_{j}X^{\mathrm{B}_{j}} \Delta^{\mathrm{A}_{i} \mathrm{~B}_{j}}\right]^{-1},
        \\
        &\text{with }\sum_{\mathrm{B}_{j}}\text{ over all sites on molecule }j\ (\mathrm{A}_{j},\mathrm{B}_{j},\mathrm{C}_{j},\ldots),
        \\
        &\text{and }\sum_{j}\text{ over all components.}
    \end{aligned}
    \label{eq:x_assoc_site}
\end{equation}

% EqID: rho_j_assoc
% Source: \cite{Chapman1990} (exact equation number not available in local PDFs)
% Status: No direct numbered source in local corpus
% Description: Provides a supporting relation used in association contribution.
% Change note: Component density relation used with association equations; no directly numbered counterpart found in local Chapman-accessible sources.

\begin{equation}
    \rho_{j}=X_{j} \rho_{\text {mixture }}
    \label{eq:rho_j_assoc}
\end{equation}

% EqID: delta_assoc
% Source: \cite{Chapman1990} (exact equation number not available in local PDFs)
% Status: No direct numbered source in local corpus
% Description: Provides a residual Helmholtz-energy relation for association contribution.
% Change note: Association strength definition used in SAFT implementations; exact numbered Chapman equation not available in local PDFs.

\begin{equation}
    \Delta^{\mathrm{A}, \mathrm{~B}_{j}}=d_{i j}{ }^3 \mathrm{g}_{i j}^{\mathrm{hs}}\left[\exp \left(\frac{\epsilon^{\mathrm{A}, \mathrm{~B}_{j}}}{k T}\right)-1\right]
    \label{eq:delta_assoc}
\end{equation}

% EqID: epsilon_assoc_mixing
% Source: \cite{Gross2002}, Eq.~(2)
% Status: Manual literature match
% Description: Provides a supporting relation used in association contribution.
% Change note: Mapped manually to the Gross 2002 cross-association energy combining rule.

\begin{equation}
    \varepsilon^{A_{i}B_{j}}=\frac{1}{2}(\varepsilon^{A_{i}B_{i}}+\varepsilon^{A_{j}B_{j}})(1-k_{ij}^{\mathrm{hb}})
    \label{eq:epsilon_assoc_mixing}
\end{equation}

% EqID: kappa_assoc_mixing
% Source: \cite{Gross2002}, Eq.~(3)
% Status: Manual literature match
% Description: Provides a supporting relation used in association contribution.
% Change note: Mapped manually to the Gross 2002 cross-association volume combining rule.

\begin{equation}
    k^{A_{i}B_{j}}=\sqrt{k^{A_{i}B_{i}}k^{A_{j}B_{j}}}\quad\left(\frac{\sqrt{\sigma_{i}\sigma_{j}}}{1/2(\sigma_{i}+\sigma_{j})}\right)^3
    \label{eq:kappa_assoc_mixing}
\end{equation}

% EqID: d_ij
% Source: \cite{Gross2001}, Eq.~(A.14)
% Status: Adapted notation
% Description: Provides a supporting relation used in association contribution.
% Change note: Mapped to the arithmetic combining rule form; this file writes it for effective diameters in association context.

\begin{equation}
    d_{i j}=\left(d_{i i}+d_{j j}\right) / 2 .
    \label{eq:d_ij}
\end{equation}

\subsection{Debye and Huckel Electrolyte Term Contribution}

See~\cite{Bulow2019}.
% EqID: ares_dh
% Source: \cite{Bulow2019}, Eq.~(2)
% Status: Adapted implementation form
% Description: Defines the Debye screening quantity used in debye and huckel electrolyte term contribution.
% Change note: Lower similarity; likely algebraically adapted for implementation or combined terms.

\begin{equation}
    \tilde{a}^{DH}=-\frac{\kappa e^{2}}{12\pi\varepsilon_{0}\varepsilon_{r}k_{B}T}\sum_{i}x_{i}z_{i}^{2}\chi_{i}
    \label{eq:ares_dh}
\end{equation}

% EqID: kappa_dh
% Source: \cite{Bulow2019}, Eq.~(2)
% Status: Adapted implementation form
% Description: Defines the Debye screening quantity used in debye and huckel electrolyte term contribution.
% Change note: Lower similarity; likely algebraically adapted for implementation or combined terms.

\begin{equation}
    \kappa=\sqrt{\frac{\rho e^{2}}{k_{B}T\varepsilon_{0}\varepsilon_{r}}\sum_{j}x_{j}z_{j}^{2}}
    \label{eq:kappa_dh}
\end{equation}

% EqID: chi_dh
% Source: \cite{Bulow2019}, Eq.~(10)
% Status: Manual literature match
% Description: Defines the Debye screening quantity used in debye and huckel electrolyte term contribution.
% Change note: Mapped manually to the ion-size function definition in the concentration-dependent dielectric derivation.

\begin{equation}
    \chi_{i}=\frac{3}{\left(\kappa d_{i}\right)^3}\left[\frac{3}{2}+\ln(1+\kappa d_{i})-2(1+\kappa d_{i})+\frac{1}{2}\left(1+\kappa d_{i}\right)^2\right]
    \label{eq:chi_dh}
\end{equation}

\subsubsection{Ion Diameter}

Temperature Independent Version
% EqID: d_ion_sigma
% Source: \cite{Bulow2019}, Eq.~(3)
% Status: Adapted implementation form
% Description: Provides a supporting relation used in debye and huckel electrolyte term contribution.
% Change note: Lower similarity; likely algebraically adapted for implementation or combined terms.

\begin{equation}
    d_{i}=\sigma_{i} \qquad i\in\mathcal{I}
    \label{eq:d_ion_sigma}
\end{equation}

Temperature Dependent Version
% EqID: d_ion_sigma_const
% Source: \cite{Bulow2019}, Eq.~(3)
% Status: Adapted implementation form
% Description: Provides a supporting relation used in debye and huckel electrolyte term contribution.
% Change note: Lower similarity; likely algebraically adapted for implementation or combined terms.

\begin{equation}
    d_{i}=\left[1-0.12\right]\sigma_{i} = .88\sigma_{i}  \qquad i\in\mathcal{I}
    \label{eq:d_ion_sigma_const}
\end{equation}

Temperature Dependent Version 2
% EqID: d_ion_barker_henderson
% Source: \cite{Bulow2019}, Eq.~(3)
% Status: Adapted implementation form
% Description: Provides a supporting relation used in debye and huckel electrolyte term contribution.
% Change note: Lower similarity; likely algebraically adapted for implementation or combined terms.

\begin{equation}
    d_{i}=\left[1-0.12 \exp \left(-3 \frac{\epsilon_{i}}{k T}\right)\right]\sigma_{i}  \qquad i\in\mathcal{I}
    \label{eq:d_ion_barker_henderson}
\end{equation}

\subsubsection{Dielectric Constant}

See~\cite{Ascani2021}.
Mole-Fraction
% EqID: epsr_mix_mole
% Source: \cite{Ascani2021}, Eq.~(3)
% Status: Adapted implementation form
% Description: Specifies dielectric-property mixing or derivative form for debye and huckel electrolyte term contribution.
% Change note: Lower similarity; likely algebraically adapted for implementation or combined terms.

\begin{equation}
    \varepsilon_{r}=\sum_{j=1}^{N_{c}}\varepsilon_{r,j}x_{j}
    \label{eq:epsr_mix_mole}
\end{equation}

Mass-Fraction
% EqID: epsr_mix_mass
% Source: \cite{Ascani2021}, Eq.~(11)
% Status: Adapted implementation form
% Description: Specifies dielectric-property mixing or derivative form for debye and huckel electrolyte term contribution.
% Change note: Lower similarity; likely algebraically adapted for implementation or combined terms.

\begin{equation}
    \varepsilon_{r}=\sum_{j=1}^{N_{c}}\varepsilon_{r,j}\,w_{j}
    =
    \frac{\sum_{j=1}^{N_{c}}\varepsilon_{r,j}\,x_{j}\,MW_{j}}{\sum_{j=1}^{N_{c}}x_{j}\,MW_{j}}
    =
    \frac{\sum_{j=1}^{N_{c}}\varepsilon_{r,j}\,x_{j}\,MW_{j}}{\overline{MW}}
    \label{eq:epsr_mix_mass}
\end{equation}

Combo
% EqID: epsr_mix_combo
% Source: \cite{Ascani2021}, Eq.~(13)
% Status: Manual literature match
% Description: Specifies dielectric-property mixing or derivative form for debye and huckel electrolyte term contribution.
% Change note: Mapped manually to the mixed solvent/ion dielectric mixing relation.

\begin{equation}
    \varepsilon_{r} = x_{sol}\,\varepsilon_{r,sol}^{(w)} + \sum_{j\in\mathcal I}\varepsilon_{r,j}\,x_{j}
    \label{eq:epsr_mix_combo}
\end{equation}

% EqID: epsr_solvent_mass
% Source: \cite{Ascani2021}, Eq.~(13)
% Status: Adapted notation
% Description: Provides a supporting relation used in debye and huckel electrolyte term contribution.
% Change note: Algebraic expansion of Eq. (13) introducing explicit solvent-only weighted dielectric form.

\begin{equation}
    \varepsilon_{r,sol}^{(w)} \equiv \sum_{j\in\mathcal S}\varepsilon_{r,j}\,w_{j}^{sol} = \frac{\sum_{j\in\mathcal S}\varepsilon_{r,j}\,x_{j}\,MW_{j}}{\sum_{j\in\mathcal S}x_{j}\,MW_{j}} = \frac{\sum_{j\in\mathcal S}\varepsilon_{r,j}\,x_{j}\,MW_{j}}{\overline{MW}_{sol}}
    \label{eq:epsr_solvent_mass}
\end{equation}

New Mixing Rule (2025)~\cite{Figiel2025}

% EqID: epsr_mix_figiel2025
% Source: \cite{Figiel2025}, Eq.~(11)
% Status: New literature extension
% Description: Provides a supporting relation used in debye and huckel electrolyte term contribution.
% Change note: Changed from earlier dielectric-mixing rules to the 2025 ion-suppression mixing form.

\begin{equation}
    \varepsilon_{r,\mathrm{mix}}=\frac{\varepsilon_{r,\mathrm{solvent,mix}}^{\mathrm{salt-free}}}{1+7.01\cdot x_{\mathrm{ion}}}
    \label{eq:epsr_mix_figiel2025}
\end{equation}

% EqID: epsr_salt_free
% Source: \cite{Figiel2025}, Eq.~(12)
% Status: New literature extension
% Description: Provides a supporting relation used in debye and huckel electrolyte term contribution.
% Change note: Changed from earlier dielectric-mixing rules to the 2025 ion-suppression mixing form.

\begin{equation}
    \varepsilon_{r,\mathrm{solvent,mix}}^{\mathrm{salt-free}}=\sum_{\mathrm{solvent}}w_{\mathrm{solvent}}^{\mathrm{salt-free}}\cdot\varepsilon_{r,\mathrm{solvent}}
    \label{eq:epsr_salt_free}
\end{equation}


Intermediates

% EqID: mw_bar
% Source: \cite{Figiel2025}, Eq.~(12) (derived helper)
% Status: Derived helper relation
% Description: Provides a supporting relation used in debye and huckel electrolyte term contribution.
% Change note: This mean-molecular-weight helper is used to evaluate Eq. (12) but is not a separately numbered equation in the paper.

\begin{equation}
    \overline{MW}=\sum_{j=1}^{N_{c}} x_{j}\,MW_{j}
    \label{eq:mw_bar}
\end{equation}

% EqID: mw_solvent_bar
% Source: \cite{Figiel2025}, Eq.~(16)
% Status: Adapted implementation form
% Description: Provides a supporting relation used in debye and huckel electrolyte term contribution.
% Change note: Lower similarity; likely algebraically adapted for implementation or combined terms.

\begin{equation}
    \overline{MW}_{sol}=\sum_{j\in\mathcal S}x_{j}\,MW_{j}
    \label{eq:mw_solvent_bar}
\end{equation}

% EqID: solvent_ion_sets
% Source: \cite{Figiel2025}, Eq.~(11) (set-definition helper)
% Status: Derived helper relation
% Description: Provides a supporting relation used in debye and huckel electrolyte term contribution.
% Change note: This solvent/ion index-set definition is implementation notation introduced in this documentation.

\begin{equation}
    \mathcal S=\{j:\,z_{j}=0\},\qquad \mathcal I=\{j:\,z_{j}\neq 0\}
    \label{eq:solvent_ion_sets}
\end{equation}

% EqID: x_solvent_total
% Source: \cite{Figiel2025}, Eq.~(16)
% Status: Adapted implementation form
% Description: Provides a supporting relation used in debye and huckel electrolyte term contribution.
% Change note: Lower similarity; likely algebraically adapted for implementation or combined terms.

\begin{equation}
    x_{sol}=\sum_{j\in\mathcal S}x_{j}
    \label{eq:x_solvent_total}
\end{equation}

\subsubsection{Bjerrum Treatment}

See~\cite{Bulow2021}.
% EqID: ares_dh_bjerrum
% Source: \cite{Bulow2021}, Eq.~(20)
% Status: Manual literature match
% Description: Defines the Debye screening quantity used in debye and huckel electrolyte term contribution.
% Change note: Mapped manually to the reduced Debye-Huckel Helmholtz form with dissociation degree factors.

\begin{equation}
    \tilde{a}^{DH}=-\frac{\kappa e^{2}}{12\pi\varepsilon_{0}\varepsilon_{r}k_{B}T}\sum_{i}\alpha_{i}x_{i}z_{i}^{2}\chi_{i}
    \label{eq:ares_dh_bjerrum}
\end{equation}

% EqID: kappa_dh_bjerrum
% Source: \cite{Bulow2021}, Eq.~(18)
% Status: Manual literature match
% Description: Defines the Debye screening quantity used in debye and huckel electrolyte term contribution.
% Change note: Mapped manually to the Bjerrum-treatment Debye screening parameter definition.

\begin{equation}
    \kappa=\sqrt{\frac{\rho e^{2}}{k_{B}T\varepsilon_{0}\varepsilon_{r}}\sum_{j}\alpha_{j}x_{j}z_{j}^{2}}
    \label{eq:kappa_dh_bjerrum}
\end{equation}

% EqID: chi_dh_bjerrum
% Source: \cite{Bulow2021}, Eq.~(19)
% Status: Manual literature match
% Description: Defines the Debye screening quantity used in debye and huckel electrolyte term contribution.
% Change note: Mapped manually to the Bjerrum-treatment chi-function definition.

\begin{equation}
    \chi_{i}=\frac{3}{\left(\kappa R_{i}\right)^3}\left[\frac{3}{2}+\ln(1+\kappa R_{i})-2(1+\kappa R_{i})+\frac{1}{2}\left(1+\kappa R_{i}\right)^2\right]
    \label{eq:chi_dh_bjerrum}
\end{equation}

\paragraph{Bjerrum Length}

% EqID: r_ion_bjerrum
% Source: \cite{Bulow2021}, Eq.~(16)
% Status: Manual literature match
% Description: Provides a supporting relation used in debye and huckel electrolyte term contribution.
% Change note: Mapped manually to the closest-approach radius selection between ion diameter and Bjerrum length.

\begin{equation}
    R_{i}=
    \begin{cases}
        d_{ion}, & d_{ion}>l_{B}, \\
        l_{B},   & d_{ion}<l_{B}.
    \end{cases}
    \label{eq:r_ion_bjerrum}
\end{equation}

% EqID: bjerrum_length
% Source: \cite{Bulow2021}, Eq.~(10)
% Status: Close literature match
% Description: Specifies dielectric-property mixing or derivative form for debye and huckel electrolyte term contribution.
% Change note: High textual similarity to a tagged equation in the cited local paper export.

\begin{equation}
    l_{B}=\frac{\left|z_{i}z_{j}\right|e^2}{8\pi\varepsilon_{0}\varepsilon_{r}k_{B}T}
    \label{eq:bjerrum_length}
\end{equation}

\paragraph{Dissociation Degree}

Solved implicitly

% EqID: alpha_ion_pair
% Source: \cite{Bulow2021}, Eq.~(14)
% Status: Close literature match
% Description: Gives an activity or fugacity-coefficient relation in debye and huckel electrolyte term contribution.
% Change note: High textual similarity to a tagged equation in the cited local paper export.

\begin{equation}
    \alpha=\frac{-1+\sqrt{1+4x_{\pm} K_{ip}\frac{\left(\gamma_{\pm}^*\left(x_{f}\right)\right)^2}{\gamma_{ip}^*}}}{2x_{\pm} K_{ip}\frac{\left(\gamma_{\pm}^*\left(x_{f}\right)\right)^2}{\gamma_{ip}^*}}
    \label{eq:alpha_ion_pair}
\end{equation}

% EqID: k_ion_pair
% Source: \cite{Bulow2021}, Eq.~(9) and Eq.~(11)
% Status: Manual literature match
% Description: Specifies dielectric-property mixing or derivative form for debye and huckel electrolyte term contribution.
% Change note: This line combines the algebraic ion-pair equilibrium form and the configurational integral expression shown as separate equations in the paper.

\begin{equation}
    K_{ip}(T)=\frac{(1-\alpha)\cdot\gamma_{ip}^{*}}{\alpha^{2}\cdot x_{\pm}\cdot\left(\gamma_{\pm}^{*}(x_{f}(\alpha))\right)^{2}} = 4\pi\rho_{N}\int_{a}^{l_{B}}\exp\left(\frac{\left|z_{i}z_{j}\right|e^2}{4\pi\varepsilon_{0}\varepsilon_{r}k_{B}T}\cdot\frac{1}{r}\right)r^2dr
    \label{eq:k_ion_pair}
\end{equation}

% EqID: gamma_ion_pair_unity
% Source: \cite{Bulow2021} (approximation not explicitly numbered)
% Status: Project-specific approximation
% Description: Provides a supporting relation used in debye and huckel electrolyte term contribution.
% Change note: Assumption $\gamma_{ip}^* \approx 1$ is used as a simplifying closure and is not a standalone numbered equation in the source paper.

\begin{equation}
    \gamma_{ip}^* \approx 1
    \label{eq:gamma_ion_pair_unity}
\end{equation}

% EqID: x_pm_balance
% Source: \cite{Bulow2021}, Eq.~(15)
% Status: Adapted notation
% Description: Provides a supporting relation used in debye and huckel electrolyte term contribution.
% Change note: Moderate-to-high similarity; notation/arrangement appears adapted from the cited equation.

\begin{equation}
    x_{\pm} = x_{f} + x_{ip}
    \label{eq:x_pm_balance}
\end{equation}

% EqID: x_free_ion_from_alpha
% Source: \cite{Bulow2021}, Eq.~(15)
% Status: Adapted implementation form
% Description: Provides a supporting relation used in debye and huckel electrolyte term contribution.
% Change note: Lower similarity; likely algebraically adapted for implementation or combined terms.

\begin{equation}
    x_{\pm} = \alpha x_{f}
    \label{eq:x_free_ion_from_alpha}
\end{equation}

% EqID: x_ion_pair_from_alpha
% Source: \cite{Bulow2021}, Eq.~(9) (derived stoichiometric form)
% Status: Derived from literature equation
% Description: Provides a supporting relation used in debye and huckel electrolyte term contribution.
% Change note: Stoichiometric rearrangement used with Eq. (9) during alpha-based ion-pair splitting.

\begin{equation}
    x_{\pm} = (1 - \alpha) x_{ip}
    \label{eq:x_ion_pair_from_alpha}
\end{equation}

% EqID: x_pm_stoich
% Source: \cite{Bulow2021}, Eq.~(7) (adapted variable form)
% Status: Derived from literature equation
% Description: Provides a supporting relation used in debye and huckel electrolyte term contribution.
% Change note: Geometric mean form mirrors the mean-ionic expression pattern and is written here for mole fractions.

\begin{equation}
    x_{\pm}=(x_{c}^{\nu c}\cdot x_{a}^{\nu a})^{\frac{1}{\nu_{c}+\nu_{a}}}
    \label{eq:x_pm_stoich}
\end{equation}

% EqID: mu_dh_infinite_dilution
% Source: \cite{Bulow2021}, Eq.~(25)
% Status: Adapted implementation form
% Description: Defines the Debye screening quantity used in debye and huckel electrolyte term contribution.
% Change note: Lower similarity; likely algebraically adapted for implementation or combined terms.

\begin{equation}
    \mu_{i}^{DH}\left(x_{f}\right)=\left(\frac{\partial A^{DH}}{\partial\rho_{i}\left(x_{f}\right)}\right)_{T,V,N_{j\neq i}}=-\frac{e^{2}z_{i}^{2}\kappa}{24\pi\varepsilon_{0}\varepsilon_{r}}\left[2\chi_{i}+\frac{\sum_{k}x_{k,f}z_{k}^{2}\sigma_{k}}{\sum_{k}x_{k,f}z_{k}^{2}}\right]
    \label{eq:mu_dh_infinite_dilution}
\end{equation}

% EqID: lngamma_i_infinite_dilution
% Source: \cite{Bulow2021}, Eq.~(13)
% Status: Manual literature match
% Description: Defines the Debye screening quantity used in debye and huckel electrolyte term contribution.
% Change note: Mapped manually to the infinite-dilution ionic activity-coefficient relation.

\begin{equation}
    \ln\gamma_{i}^{*}\left(x_{f,i}\right)=-\frac{e^{2}z_{i}^{2}\kappa}{24\pi\varepsilon_{0}\varepsilon_{r}}\left[2\chi_{i}+\frac{\sum_{k}x_{k,f}z_{k}^{2}\sigma_{k}}{\sum_{k}x_{k,f}z_{k}^{2}}\right]
    \label{eq:lngamma_i_infinite_dilution}
\end{equation}

% EqID: gamma_pm_x
% Source: \cite{Bulow2021}, Eq.~(26)
% Status: Adapted implementation form
% Description: Gives an activity or fugacity-coefficient relation in debye and huckel electrolyte term contribution.
% Change note: Lower similarity; likely algebraically adapted for implementation or combined terms.

\begin{equation}
    \gamma_{\pm}^{*}=(\gamma_{c}^{*,\nu c}\cdot\gamma_{a}^{*,\nu a})^{\frac{1}{\nu_{c}+\nu_{a}}}
    \label{eq:gamma_pm_x}
\end{equation}

\subsection{Born Electrolyte Term Contribution}

See~\cite{Bulow2021a}.
% EqID: ares_born
% Source: \cite{Bulow2021a}, Eq.~(2)
% Status: Adapted notation
% Description: Provides a residual Helmholtz-energy relation for born electrolyte term contribution.
% Change note: Moderate-to-high similarity; notation/arrangement appears adapted from the cited equation.

\begin{equation}
    \tilde{a}^{\text {Born }}(\varepsilon(x))=-\frac{e^2}{4 \pi \varepsilon_{0} k_{B} T}\left(1-\frac{1}{\varepsilon_{r,\mathrm{bulk}}}\right) \sum_{i} \frac{x_{i} z_{i}^2}{d^{\text{Born}}_{i}}
    \label{eq:ares_born}
\end{equation}

% EqID: d_born_equals_d
% Source: \cite{Bulow2021a} (modeling assumption, not separately numbered)
% Status: No direct numbered source in local corpus
% Description: Provides a supporting relation used in born electrolyte term contribution.
% Change note: The equality $d_i^{\mathrm{Born}}=d_i$ is a modeling assumption used in implementations; no standalone numbered equation in Part I was found.

\begin{equation}
    d^{\text{Born}}_{i} = d_{i} \qquad i\in\mathcal{I}
    \label{eq:d_born_equals_d}
\end{equation}

\subsubsection{Solvation Shell Model + Dielectric Saturation}

Version 1a: Equation taken directly from the paper (probably wrong):~\cite{Figiel2025}

% EqID: ares_born_ssmds_v1
% Source: \cite{Figiel2025}, Eq.~(7)
% Status: Modified/flagged form
% Description: Provides a supporting relation used in born electrolyte term contribution.
% Change note: Changed from literature-transcribed form; note says this direct paper form is probably wrong.

\begin{equation}
    \begin{aligned}
        a_{\mathrm{SSM+DS}}^{\mathrm{Born}}
        &=-\frac{e^{2}}{4\pi\varepsilon_{0}k_{\mathrm{B}}T}
        \left(1-\frac{1}{\varepsilon_{r,\mathrm{bulk}}}\right)
        \sum_{i}\frac{x_{i}z_{i}^{2}}{d_{i}^{\mathrm{Born}}+\Delta d_{i}}
        \\
        &\quad +\left(1-\frac{1}{\varepsilon_{r,\mathrm{bulk}}}\right)
        \left(\sum_{i}x_{i}z_{i}^{2}
        \left(\frac{1}{d_{i}^{\mathrm{Born}}}-\frac{1}{d_{i}^{\mathrm{Born}}+\Delta d_{i}}\right)\right)
    \end{aligned}
    \label{eq:ares_born_ssmds_v1}
\end{equation}

Version 1b: Equation from paper with modification to include the first fraction for the DS term (probably wrong):~\cite{Figiel2025}

% EqID: ares_born_ssmds_v2
% Source: \cite{Figiel2025}, Eq.~(7)
% Status: Modified form
% Description: Provides a supporting relation used in born electrolyte term contribution.
% Change note: Changed from Version 1a by adding the first DS prefactor term (documented in-note).

\begin{equation}
    \begin{aligned}
        a_{\mathrm{SSM+DS}}^{\mathrm{Born}}
        &=-\frac{e^{2}}{4\pi\varepsilon_{0}k_{\mathrm{B}}T}
        \left(1-\frac{1}{\varepsilon_{r,\mathrm{bulk}}}\right)
        \sum_{i}\frac{x_{i}z_{i}^{2}}{d_{i}^{\mathrm{Born}}+\Delta d_{i}}
        \\
        &\quad -\frac{e^{2}}{4\pi\varepsilon_{0}k_{\mathrm{B}}T}
        \left(1-\frac{1}{\varepsilon_{r,\mathrm{bulk}}}\right)
        \left(\sum_{i}x_{i}z_{i}^{2}
        \left(\frac{1}{d_{i}^{\mathrm{Born}}}-\frac{1}{d_{i}^{\mathrm{Born}}+\Delta d_{i}}\right)\right)
    \end{aligned}
    \label{eq:ares_born_ssmds_v2}
\end{equation}

Version 2 (my version): potential fix (changing the second $\varepsilon_{r,\mathrm{solv}}$ to $\varepsilon_{r,\mathrm{ion}}$ and including the first fraction for the DS term)

% EqID: ares_born_ssm
% Source: \cite{Figiel2025}, Eq.~(7)
% Status: Project-specific modification
% Description: Provides a supporting relation used in born electrolyte term contribution.
% Change note: Changed from paper-style form by switching to ion dielectric in DS term and splitting SSM/DS terms.

\begin{equation}
    a_{\mathrm{SSM}}^{\mathrm{Born}}=-\frac{e^{2}}{4\pi\varepsilon_{0}k_{\mathrm{B}}T}\sum_{i}x_{i}z_{i}^{2} \left(1-\frac{1}{\varepsilon_{r,\mathrm{bulk}}}\right)\left(\frac{1}{d_{i}^{\mathrm{Born}}+\Delta d_{i}}\right)
    \label{eq:ares_born_ssm}
\end{equation}

% EqID: ares_born_ds
% Source: \cite{Figiel2025}, Eq.~(7)
% Status: Project-specific modification
% Description: Provides a supporting relation used in born electrolyte term contribution.
% Change note: Changed from paper-style form by switching to ion dielectric in DS term and splitting SSM/DS terms.

\begin{equation}
    a_{\mathrm{DS}}^{\mathrm{Born}}=-\frac{e^{2}}{4\pi\varepsilon_{0}k_{\mathrm{B}}T}\sum_{i}x_{i}z_{i}^{2} \left(1-\frac{1}{\varepsilon_{r,\mathrm{ion}}}\right)\left(\frac{1}{d_{i}^{\mathrm{Born}}}-\frac{1}{d_{i}^{\mathrm{Born}}+\Delta d_{i}}\right)
    \label{eq:ares_born_ds}
\end{equation}

% EqID: ares_born_ssmds_v3
% Source: \cite{Figiel2025}, Eq.~(7)
% Status: Adapted notation
% Description: Provides a supporting relation used in born electrolyte term contribution.
% Change note: Moderate-to-high similarity; notation/arrangement appears adapted from the cited equation.

\begin{equation}
    \begin{aligned}
        a_{\mathrm{SSM+DS}}^{\mathrm{Born}}
        &=-\frac{e^{2}}{4\pi\varepsilon_{0}k_{\mathrm{B}}T}
        \sum_{i}x_{i}z_{i}^{2}
        \left(1-\frac{1}{\varepsilon_{r,\mathrm{bulk}}}\right)
        \left(\frac{1}{d_{i}^{\mathrm{Born}}+\Delta d_{i}}\right)
        \\
        &\quad -\frac{e^{2}}{4\pi\varepsilon_{0}k_{\mathrm{B}}T}
        \sum_{i}x_{i}z_{i}^{2}
        \left(1-\frac{1}{\varepsilon_{r,\mathrm{ion}}}\right)
        \left(\frac{1}{d_{i}^{\mathrm{Born}}}-\frac{1}{d_{i}^{\mathrm{Born}}+\Delta d_{i}}\right)
    \end{aligned}
    \label{eq:ares_born_ssmds_v3}
\end{equation}


% EqID: delta_d_born
% Source: \cite{Figiel2025}, Eq.~(8)
% Status: Close literature match
% Description: Provides a supporting relation used in born electrolyte term contribution.
% Change note: High textual similarity to a tagged equation in the cited local paper export.

\begin{equation}
    \Delta d_{i}=\frac{(f_{mix}-1)}{|z_{i}|}\cdot d_{i}^{\mathrm{Born}}
    \label{eq:delta_d_born}
\end{equation}

% EqID: f_mix
% Source: \cite{Figiel2025}, Eq.~(10)
% Status: Close literature match
% Description: Provides a supporting relation used in born electrolyte term contribution.
% Change note: High textual similarity to a tagged equation in the cited local paper export.

\begin{equation}
    f_{mix}=\sum_{k}x_{k}f_{k}
    \label{eq:f_mix}
\end{equation}

% EqID: d_born_fitted
% Source: \cite{Figiel2025}, Eq.~(6)-Eq.~(8) (parameterization choice)
% Status: No direct numbered source in local corpus
% Description: Provides a supporting relation used in born electrolyte term contribution.
% Change note: Using fitted Born diameters is discussed as a parameterization choice; the exact identity line here is not separately numbered.

\begin{equation}
    d^{\text{Born}}_{i} = d^{\text{Born, fitted}}_{i} \qquad i\in\mathcal{I}
    \label{eq:d_born_fitted}
\end{equation}

\section{Pressure (Compressibility Factor)}

See~\cite{Gross2001}.
% EqID: pressure_from_z
% Source: \cite{Gross2001}, Eq.~(A.21)
% Status: Adapted implementation form
% Description: Provides a supporting relation used in pressure (compressibility factor).
% Change note: Lower similarity; likely algebraically adapted for implementation or combined terms.

\begin{equation}
    P=ZkT\rho\left(10^{10}\frac{\hat{\mathrm{A}}}{\mathrm{m}}\right)^{3}
    \label{eq:pressure_from_z}
\end{equation}

% EqID: z_from_eta
% Source: \cite{Gross2001}, Eq.~(A.22)
% Status: Close literature match
% Description: Provides a differential relation needed for pressure (compressibility factor) calculations.
% Change note: High textual similarity to a tagged equation in the cited local paper export.

\begin{equation}
    Z=1+\eta\left(\frac{\partial\tilde{a}^\mathrm{res}}{\partial\eta}\right)_{T,x_{i}}
    \label{eq:z_from_eta}
\end{equation}
% EqID: z_eta_rho_identity
% Source: \cite{Gross2001}, Eq.~(A.22)
% Status: Adapted implementation form
% Description: Provides a differential relation needed for pressure (compressibility factor) calculations.
% Change note: Lower similarity; likely algebraically adapted for implementation or combined terms.

\begin{equation}
    \eta\left(\frac{\partial\tilde{a}^\mathrm{res}}{\partial\eta}\right)_{T,x_{i}} = \rho\left(\frac{\partial\tilde{a}^\mathrm{res}}{\partial\rho}\right)_{T,x_{i}}
    \label{eq:z_eta_rho_identity}
\end{equation}

% EqID: z_from_rho
% Source: \cite{Gross2001}, Eq.~(A.22)
% Status: Close literature match
% Description: Provides a differential relation needed for pressure (compressibility factor) calculations.
% Change note: High textual similarity to a tagged equation in the cited local paper export.

\begin{equation}
    Z=1+\rho\left(\frac{\partial\tilde{a}^\mathrm{res}}{\partial\rho}\right)_{T,x_{i}}
    \label{eq:z_from_rho}
\end{equation}

% EqID: z_total
% Source: \cite{Gross2001}, Eq.~(A.24)
% Status: Project-specific extension
% Description: Provides a supporting relation used in pressure (compressibility factor).
% Change note: Mapped to Eq. (A.24) and extended here by adding association, Debye-Huckel, and Born compressibility contributions.

\begin{equation}
    Z = 1 + Z^{hc} + Z^{disp} + Z^{assoc} + Z^{DH} + Z^{Born}
    \label{eq:z_total}
\end{equation}

% EqID: z_minus_one_sum
% Source: \cite{Gross2001}, Eq.~(A.24) with the documented ePC-SAFT term decomposition
% Status: Project-specific extension
% Description: Gives the explicit residual compressibility identity in pressure (compressibility factor).
% Change note: Added as the direct $Z-1$ form so later fugacity equations can use only $Z$ and $Z^\alpha$ terms explicitly.

\begin{equation}
    Z - 1 = Z^{hc} + Z^{disp} + Z^{assoc} + Z^{DH} + Z^{Born}
    \label{eq:z_minus_one_sum}
\end{equation}

% EqID: z_alpha
% Source: \cite{Gross2001}, Eq.~(A.22)
% Status: Adapted notation
% Description: Provides a differential relation needed for pressure (compressibility factor) calculations.
% Change note: Moderate-to-high similarity; notation/arrangement appears adapted from the cited equation.

\begin{equation}
    Z^\alpha =\rho\left(\frac{\partial\tilde{a}^\alpha}{\partial\rho}\right)_{T,x}
    .
    \label{eq:z_alpha}
\end{equation}

\subsection{Hard-Chain Reference Contribution}

See~\cite{Gross2001}.
% EqID: z_hc
% Source: \cite{Gross2001}, Eq.~(A.22)
% Status: Adapted notation
% Description: Provides a differential relation needed for hard-chain reference contribution calculations.
% Change note: Moderate-to-high similarity; notation/arrangement appears adapted from the cited equation.

\begin{equation}
    Z^{hc} =\rho\left(\frac{\partial\tilde{a}^{hc}}{\partial\rho}\right)_{T,x}
    .
    \label{eq:z_hc}
\end{equation}

% EqID: z_hc_explicit
% Source: \cite{Gross2001}, Eq.~(A.25)
% Status: Close literature match
% Description: Provides a differential relation needed for hard-chain reference contribution calculations.
% Change note: High textual similarity to a tagged equation in the cited local paper export.

\begin{equation}
    Z^\mathrm{hc}=\bar{m}Z^\mathrm{hs}-\sum_{i}x_{\mathrm{i}}(m_{i}-1)(g_{ii}^\mathrm{hs})^{-1}\rho\frac{\partial g_{ii}^\mathrm{hs}}{\partial\rho}
    \label{eq:z_hc_explicit}
\end{equation}

% EqID: z_hs
% Source: \cite{Gross2001}, Eq.~(A.26)
% Status: Manual literature match
% Description: Provides a differential relation needed for hard-chain reference contribution calculations.
% Change note: Mapped manually to the hard-sphere compressibility expression.

\begin{equation}
    Z^{hs} =\rho\left(\frac{\partial\tilde{a}^{hs}}{\partial\rho}\right)_{T,x}
    .
    \label{eq:z_hs}
\end{equation}

% EqID: z_hs_explicit
% Source: \cite{Gross2001}, Eq.~(A.26)
% Status: Manual literature match
% Description: Provides a supporting relation used in hard-chain reference contribution.
% Change note: Mapped manually to the explicit hard-sphere compressibility formula.

\begin{equation}
    Z^{\mathrm{hs}}=\frac{\zeta_{3}}{(1-\zeta_{3})}+\frac{3\zeta_{1}\zeta_{2}}{\zeta_{0}(1-\zeta_{3})^{2}}+\frac{3\zeta_{2}^{3}-\zeta_{3}\zeta_{2}^{3}}{\zeta_{0}(1-\zeta_{3})^{3}}
    \label{eq:z_hs_explicit}
\end{equation}

% EqID: dg_hs_drho
% Source: \cite{Gross2001}, Eq.~(A.27)
% Status: Manual literature match
% Description: Provides a differential relation needed for hard-chain reference contribution calculations.
% Change note: Mapped manually to the density derivative of the contact value relation.

\begin{equation}
    \begin{aligned}
        \rho\frac{\partial g_{ij}^{\mathrm{hs}}}{\partial\rho}=\frac{\zeta_{3}}{\left(1-\zeta_{3}\right)^{2}}+\left(\frac{d_{i}d_{j}}{d_{i}+d_{j}}\right)\left(\frac{3\zeta_{2}}{\left(1-\zeta_{3}\right)^{2}}+\frac{6\zeta_{2}\zeta_{3}}{\left(1-\zeta_{3}\right)^{3}}\right) \\+\left(\frac{d_{i}d_{j}}{d_{i}+d_{j}}\right)^{2}\left(\frac{4{\zeta_{2}}^{2}}{\left(1-{\zeta_{3}}\right)^{3}}+\frac{6{\zeta_{2}}^{2}\zeta_{3}}{\left(1-{\zeta_{3}}\right)^{4}}\right)
    \end{aligned}
    \label{eq:dg_hs_drho}
\end{equation}

\subsection{Dispersion Contribution}

See~\cite{Gross2001}.
% EqID: z_disp
% Source: \cite{Gross2001}, Eq.~(A.28)
% Status: Manual literature match
% Description: Provides a differential relation needed for dispersion contribution calculations.
% Change note: Mapped manually to the dispersion compressibility contribution.

\begin{equation}
    Z^{disp} =\rho\left(\frac{\partial\tilde{a}^{disp}}{\partial\rho}\right)_{T,x}
    .
    \label{eq:z_disp}
\end{equation}

% EqID: z_disp_explicit
% Source: \cite{Gross2001}, Eq.~(A.28)
% Status: Close literature match
% Description: Provides a differential relation needed for dispersion contribution calculations.
% Change note: High textual similarity to a tagged equation in the cited local paper export.

\begin{equation}
    \begin{aligned}
        Z^{\mathrm{disp}}
        &=-2\pi\rho\frac{\partial(\eta I_{1})}{\partial\eta}\overline{m^{2}\epsilon\sigma^{3}}
        \\
        &\quad -\pi\rho\bar{m}\left[C_{1}\frac{\partial(\eta I_{2})}{\partial\eta}+C_{2}\eta I_{2}\right]\overline{m^{2}\epsilon^{2}\sigma^{3}}
    \end{aligned}
    \label{eq:z_disp_explicit}
\end{equation}

% EqID: deta_i1_deta
% Source: \cite{Gross2001}, Eq.~(A.30)
% Status: Adapted implementation form
% Description: Provides a differential relation needed for dispersion contribution calculations.
% Change note: Lower similarity; likely algebraically adapted for implementation or combined terms.

\begin{equation}
    \frac{\partial(\eta I_{1})}{\partial\eta}=\sum_{j=0}^6a_{j}(\bar{m})(j+1)\eta^j
    \label{eq:deta_i1_deta}
\end{equation}

% EqID: deta_i2_deta
% Source: \cite{Gross2001}, Eq.~(A.30)
% Status: Adapted implementation form
% Description: Provides a differential relation needed for dispersion contribution calculations.
% Change note: Lower similarity; likely algebraically adapted for implementation or combined terms.

\begin{equation}
    \frac{\partial(\eta I_{2})}{\partial\eta}=\sum_{i=0}^6b_{j}(\bar{m})(j+1)\eta^i
    \label{eq:deta_i2_deta}
\end{equation}

% EqID: c2_disp
% Source: \cite{Gross2001}, Eq.~(A.31)
% Status: Close literature match
% Description: Provides a differential relation needed for dispersion contribution calculations.
% Change note: High textual similarity to a tagged equation in the cited local paper export.

\begin{equation}
    C_{2}
    =\frac{\partial C_{1}}{\partial\eta}
    =- C_{1}^{2}\biggl(
    \bar{m}\frac{-4\eta^{2}+20\eta+8}{\left(1-\eta\right)^{5}}
    +(1-\bar{m})\frac{2\eta^{3}+12\eta^{2}-48\eta+40}{\left[(1-\eta)(2-\eta)\right]^{3}}
    \biggr)
    \label{eq:c2_disp}
\end{equation}

\subsection{Association Contribution}

See~\cite{Chapman1990}.
% EqID: z_assoc
% Source: \cite{Chapman1990} (exact equation number not available in local PDFs)
% Status: No direct numbered source in local corpus
% Description: Provides a differential relation needed for association contribution calculations.
% Change note: Association compressibility relation is a derivative form from SAFT association theory; exact numbered source is unavailable in local PDFs.

\begin{equation}
    Z^{assoc} =\rho\left(\frac{\partial\tilde{a}^{assoc}}{\partial\rho}\right)_{T,x}
    .
    \label{eq:z_assoc}
\end{equation}

% EqID: z_assoc_explicit
% Source: \cite{Chapman1990} (exact equation number not available in local PDFs)
% Status: No direct numbered source in local corpus
% Description: Provides a differential relation needed for association contribution calculations.
% Change note: Expanded association compressibility expression; no direct numbered counterpart found in local Chapman-accessible files.

\begin{equation}
    Z^{assoc}=\rho\sum_{i}x_{i}\sum_{A_{i}}\left(\frac{1}{X^{A_{i}}}-\frac{1}{2}\right)\left(\frac{\partial X^{A_{i}}}{\partial\rho}\right)_{T,x}.
    \label{eq:z_assoc_explicit}
\end{equation}

% EqID: dx_assoc_drho
% Source: \cite{Chapman1990} (exact equation number not available in local PDFs)
% Status: No direct numbered source in local corpus
% Description: Provides a differential relation needed for association contribution calculations.
% Change note: Site-fraction density derivative closure used in association compressibility calculations; exact numbered reference unavailable locally.

\begin{equation}
    \begin{aligned}
        \left(\frac{\partial X^{A_{i}}}{\partial\rho}\right)_{T,x}
        &=-(X^{A_{i}})^{2}\sum_{j}\sum_{B_{j}}x_{j}
        \Bigg[
        X^{B_{j}}\Delta^{A_{i}B_{j}}
        \\
        &\quad +\rho\Bigg(
        \Delta^{A_{i}B_{j}}\left(\frac{\partial X^{B_{j}}}{\partial\rho}\right)_{T,x}
        +X^{B_{j}}\left(\frac{\partial\Delta^{A_{i}B_{j}}}{\partial\rho}\right)_{T,x}
        \Bigg)
        \Bigg]
    \end{aligned}
    \label{eq:dx_assoc_drho}
\end{equation}

% EqID: ddelta_assoc_drho
% Source: \cite{Chapman1990} (exact equation number not available in local PDFs)
% Status: No direct numbered source in local corpus
% Description: Defines the Debye screening quantity used in association contribution.
% Change note: Derivative of association strength term used in implementation; no direct numbered source in local Chapman files.

\begin{equation}
    \left(\frac{\partial\Delta^{A_{i}B_{j}}}{\partial\rho}\right)_{T,x}=d_{ij}^3\kappa^{A_{i}B_{j}}\left[\exp(\epsilon^{A_{i}B_{j}}/kT)-1\right]\left(\frac{\partial g_{ij}^{hs}}{\partial\rho}\right)_{T,x}
    \label{eq:ddelta_assoc_drho}
\end{equation}


\subsection{Debye and Huckel Electrolyte Term Contribution}

% EqID: z_dh
% Source: \cite{Cameretti2005}, Eq.~(18)
% Status: Manual literature match
% Description: Provides a differential relation needed for debye and huckel electrolyte term contribution calculations.
% Change note: Mapped manually to the Debye-Huckel pressure contribution form.

\begin{equation}
    Z^{DH} =\rho\left(\frac{\partial\tilde{a}^{DH}}{\partial\rho}\right)_{T,x}
    .
    \label{eq:z_dh}
\end{equation}

See~\cite{Cameretti2005}.
% EqID: z_dh_explicit
% Source: \cite{Cameretti2005}, Eq.~(10)
% Status: Adapted implementation form
% Description: Defines the Debye screening quantity used in debye and huckel electrolyte term contribution.
% Change note: Lower similarity; likely algebraically adapted for implementation or combined terms.

\begin{equation}
    Z^{DH}=-\frac{\kappa e^2}{24\pi kT\epsilon}\sum_{i}x_{i}z_{i}{}^{2}\sigma_{i}
    \label{eq:z_dh_explicit}
\end{equation}

% EqID: sigma_dh
% Source: \cite{Cameretti2005}, Eq.~(20)
% Status: Close literature match
% Description: Defines the Debye screening quantity used in debye and huckel electrolyte term contribution.
% Change note: High textual similarity to a tagged equation in the cited local paper export.

\begin{equation}
    \sigma_{i}=\left(\frac{\partial(\kappa\chi_{i})}{\partial\kappa}\right)_{T,\mathrm{x}}=-2\chi_{i}+\frac{3}{1+\kappa a_{i}}
    \label{eq:sigma_dh}
\end{equation}

\subsubsection{Bjerrum Treatment}

See~\cite{Bulow2021}.
% EqID: z_dh_bjerrum
% Source: \cite{Bulow2021}, Eq.~(18)
% Status: Adapted implementation form
% Description: Defines the Debye screening quantity used in debye and huckel electrolyte term contribution.
% Change note: Lower similarity; likely algebraically adapted for implementation or combined terms.

\begin{equation}
    Z^{DH}=-\frac{\kappa e^2}{24\pi kT\epsilon}\sum_{i}\alpha _{i}x_{i}z_{i}{}^{2}\sigma_{i}
    \label{eq:z_dh_bjerrum}
\end{equation}

% EqID: sigma_dh_bjerrum
% Source: \cite{Bulow2021}, Eq.~(26)
% Status: Adapted implementation form
% Description: Defines the Debye screening quantity used in debye and huckel electrolyte term contribution.
% Change note: Lower similarity; likely algebraically adapted for implementation or combined terms.

\begin{equation}
    \sigma_{i}=\left(\frac{\partial(\kappa\chi_{i})}{\partial\kappa}\right)_{T,\mathrm{x}}=-2\chi_{i}+\frac{3}{1+\kappa R_{i}}
    \label{eq:sigma_dh_bjerrum}
\end{equation}

\subsection{Born Electrolyte Term Contribution}

% EqID: z_born
% Source: \cite{Bulow2021a}, Eq.~(4)
% Status: Manual literature match
% Description: Provides a differential relation needed for born electrolyte term contribution calculations.
% Change note: Mapped manually to the Born-pressure-zero relation; this compressibility form is the equivalent EOS representation.

\begin{equation}
    Z^{Born} =\rho\left(\frac{\partial\tilde{a}^{Born}}{\partial\rho}\right)_{T,x}
    .
    \label{eq:z_born}
\end{equation}

% EqID: z_born_zero
% Source: \cite{Bulow2021a}, Eq.~(4)
% Status: Manual literature match
% Description: Provides a supporting relation used in born electrolyte term contribution.
% Change note: Direct zero Born-pressure statement from the Part I formulation.

\begin{equation}
    Z^{Born}= 0
    \label{eq:z_born_zero}
\end{equation}



\section{Chemical Potential}

See~\cite{Gross2001}.
% EqID: mu_res
% Source: \cite{Gross2001}, Eq.~(A.1)
% Status: Adapted implementation form
% Description: Gives a chemical-potential contribution in chemical potential.
% Change note: Lower similarity; likely algebraically adapted for implementation or combined terms.

\begin{equation}
    \tilde{\mu_{k}}^{\mathrm{res}}=\frac{\mu_{k}^\mathrm{res}}{kT} = \frac{\hat{\mu}_{k}^\mathrm{res}}{RT}
    \label{eq:mu_res}
\end{equation}

% EqID: mu_res_from_ares
% Source: \cite{Gross2001}, Eq.~(A.33)
% Status: Close literature match
% Description: Gives a chemical-potential contribution in chemical potential.
% Change note: High textual similarity to a tagged equation in the cited local paper export.

\begin{equation}
    \begin{aligned}
        \tilde{\mu_{k}}^{\mathrm{res}}
        &=\tilde{a}^{\mathrm{res}}+(\mathrm{Z}-1)
        +\left(\frac{\partial\tilde{a}^{\mathrm{res}}}{\partial x_{k}}\right)_{T,\nu,x_{i\neq k}}
        \\
        &\quad -\sum_{j=1}^{N}\left[x_{j}\left(\frac{\partial\tilde{a}^{\mathrm{res}}}{\partial x_{j}}\right)_{T,\nu,x_{i\neq j}}\right]
    \end{aligned}
    \label{eq:mu_res_from_ares}
\end{equation}

% EqID: mu_res_sum
% Source: \cite{Gross2001}, Eq.~(A.33)
% Status: Manual literature match
% Description: Gives the explicit residual chemical-potential decomposition in chemical potential.
% Change note: Written in explicit non-summation form to match the style used for the other property families in this section.

\begin{equation}
    \tilde{\mu_{k}}^{\mathrm{res}}=\tilde{\mu_{k}}^{\mathrm{hc}}+\tilde{\mu_{k}}^{\mathrm{disp}}+\tilde{\mu_{k}}^{\mathrm{assoc}}+\tilde{\mu_{k}}^{\mathrm{DH}}+\tilde{\mu_{k}}^{\mathrm{Born}}
    \label{eq:mu_res_sum}
\end{equation}

% EqID: zeta_n_xk
% Source: \cite{Gross2001}, Eq.~(A.34)
% Status: Close literature match
% Description: Provides a differential relation needed for chemical-potential calculations.
% Change note: High textual similarity to a tagged equation in the cited local paper export.

\begin{equation}
    \zeta_{n,xk}=\left(\frac{\partial\zeta_{n}}{\partial x_{k}}\right)_{T,\boldsymbol{\rho},x_{j\neq k}}=\frac{\pi}{6}\rho m_{k}(d_{k})^{n}\quad n\in\{0,1,2,3\}
    \label{eq:zeta_n_xk}
\end{equation}

\subsection{Hard-Chain Reference Contribution}

See~\cite{Gross2001}.
% EqID: mu_hc
% Source: \cite{Gross2001}, Eq.~(A.33)
% Status: Manual literature match
% Description: Gives a chemical-potential contribution in hard-chain reference contribution.
% Change note: Eq. (A.33) specialization to the hard-chain contribution.

\begin{equation}
    \begin{aligned}
        \tilde{\mu}^{hc}_{k}
        &=\tilde{a}^{\mathrm{hc}} + Z^{hc}
        +\left(\frac{\partial\tilde{a}^{\mathrm{hc}}}{\partial x_{k}}\right)_{T,\nu,x_{i\neq k}}
        \\
        &\quad -\sum_{j=1}^{N}\left[x_{j}\left(\frac{\partial\tilde{a}^{\mathrm{hc}}}{\partial x_{j}}\right)_{T,\nu,x_{i\neq j}}\right]
    \end{aligned}
    \label{eq:mu_hc}
\end{equation}

% EqID: dares_hc_dxk
% Source: \cite{Gross2001}, Eq.~(A.35) (appendix sequence inference)
% Status: Manual literature match
% Description: Provides a differential relation needed for hard-chain reference contribution calculations.
% Change note: Mapped by appendix sequence between Eq. (A.34) and Eq. (A.36); this line corresponds to the hard-sphere composition derivative expression.

\begin{equation}
    \begin{aligned}
        \left(\frac{\partial\tilde{a}^{\mathrm{hc}}}{\partial x_{k}}\right)_{T,\boldsymbol{\rho},x_{j\neq k}}
        &=m_{k}\tilde{a}^{\mathrm{hs}}
        +\bar{m}\left(\frac{\partial\tilde{a}^{\mathrm{hs}}}{\partial x_{k}}\right)_{T,\boldsymbol{\rho},x_{j\neq k}}
        \\
        &\quad -\sum_{i}x_{i}(m_{i}-1)(g_{ii}^{\mathrm{hs}})^{-1}
        \left(\frac{\partial g_{ii}^{\mathrm{hs}}}{\partial x_{k}}\right)_{T,\boldsymbol{\rho},x_{j\neq k}}
    \end{aligned}
    \label{eq:dares_hc_dxk}
\end{equation}

% EqID: dares_hs_dxk
% Source: \cite{Gross2001}, Eq.~(A.36)
% Status: Manual literature match
% Description: Provides a differential relation needed for hard-chain reference contribution calculations.
% Change note: Mapped manually to the hard-sphere composition derivative expression.

\begin{equation}
    \begin{aligned}
        \left(\frac{\partial\tilde{a}^{\mathrm{hs}}}{\partial x_{k}}\right)_{T,\boldsymbol{\rho},x_{j\neq k}}
        &=- \frac{\zeta_{0,xk}}{\zeta_{0}}\tilde{a}^{\mathrm{hs}}+\frac{1}{\zeta_{0}}\biggl[
        \frac{3(\zeta_{1,xk}\zeta_{2}+\zeta_{1}\zeta_{2,xk})}{(1-\zeta_{3})}
        +\frac{3\zeta_{1}\zeta_{2}\zeta_{3,xk}}{\left(1-\zeta_{3}\right)^{2}}
        \\
        &\quad +\frac{3\zeta_{2}^{2}\zeta_{2,xk}}{\zeta_{3}(1-\zeta_{3})^{2}}
        +\frac{\zeta_{2}^{3}\zeta_{3,xk}(3\zeta_{3}-1)}{\zeta_{3}^{2}(1-\zeta_{3})^{3}}
        \\
        &\quad +\left(
        \frac{3{\zeta_{2}}^{2}{\zeta_{2,xk}}{\zeta_{3}}-2{\zeta_{2}}^{3}{\zeta_{3,xk}}}{{\zeta_{3}}^{3}}
        -{\zeta_{0,xk}}
        \right)\ln(1-{\zeta_{3}})
        \\
        &\quad +\left(\zeta_{0}-\frac{\zeta_{2}^{3}}{\zeta_{3}^{2}}\right)\frac{\zeta_{3,xk}}{(1-\zeta_{3})}
        \biggr]
    \end{aligned}
    \label{eq:dares_hs_dxk}
\end{equation}

% EqID: dg_hs_dxk
% Source: \cite{Gross2001}, Eq.~(A.37)
% Status: Close literature match
% Description: Provides a differential relation needed for hard-chain reference contribution calculations.
% Change note: High textual similarity to a tagged equation in the cited local paper export.

\begin{equation}
    \begin{aligned}
        \left(\frac{\partial g_{ij}^{\mathrm{hs}}}{\partial x_{k}}\right)_{T,\boldsymbol{\rho},\boldsymbol{\rho}_{j\neq k}}
        &=\frac{\zeta_{3,x,k}}{\left(1-\zeta_{3}\right)^{2}}
        +\left(\frac{d_{i}d_{j}}{d_{i}+d_{j}}\right)
        \left(\frac{3\zeta_{2,x,k}}{(1-\zeta_{3})^{2}}+\frac{6\zeta_{2}\zeta_{3,x,k}}{(1-\zeta_{3})^{3}}\right)
        \\
        &\quad +\left(\frac{d_{i}d_{j}}{d_{i}+d_{j}}\right)^{2}
        \left(\frac{4\zeta_{2}\zeta_{2,x,k}}{(1-\zeta_{3})^{3}}+\frac{6\zeta_{2}^{2}\zeta_{3,x,k}}{(1-\zeta_{3})^{4}}\right)
    \end{aligned}
    \label{eq:dg_hs_dxk}
\end{equation}

\subsection{Dispersion Contribution}

See~\cite{Gross2001}.
% EqID: mu_disp
% Source: \cite{Gross2001}, Eq.~(A.33)
% Status: Manual literature match
% Description: Gives a chemical-potential contribution in dispersion contribution.
% Change note: Eq. (A.33) specialization to the dispersion contribution.

\begin{equation}
    \begin{aligned}
        \tilde{\mu}^{disp}_{k}
        &=\tilde{a}^{\mathrm{disp}} + Z^{disp}
        +\left(\frac{\partial\tilde{a}^{\mathrm{disp}}}{\partial x_{k}}\right)_{T,\nu,x_{i\neq k}}
        \\
        &\quad -\sum_{j=1}^{N}\left[x_{j}\left(\frac{\partial\tilde{a}^{\mathrm{disp}}}{\partial x_{j}}\right)_{T,\nu,x_{i\neq j}}\right]
    \end{aligned}
    \label{eq:mu_disp}
\end{equation}

\begin{equation}
    \begin{aligned}
        \left(\frac{\partial \tilde{a}^{\mathrm{disp}}}{\partial x_k}\right)_{T, \boldsymbol{p}, x_{j \neq k}}
        =& -2 \pi \rho\left[I_{1, x k} \overline{m^2 \epsilon \sigma^3}+I_1 \overline{\left(m^2 \epsilon \sigma^3\right)_{x k}}\right] \\
        &- \pi \rho\left\{\left[m_k C_1 I_2+\bar{m} C_{1, x k} I_2+\bar{m} C_1 I_{2, x k}\right] \overline{m^2 \epsilon^2 \sigma^3}+\right.
        \left.\bar{m} C_1 I_2\left(\bar{m}^2 \epsilon^2 \sigma^3\right)_{x k}\right\}
    \end{aligned}
\end{equation}

% EqID: m2epssigma3_xk
% Source: \cite{Gross2001}, Eq.~(A.39) (appendix sequence inference)
% Status: Manual literature match
% Description: Provides a residual Helmholtz-energy relation for dispersion contribution.
% Change note: Mapped by appendix sequence as the first dispersion composition derivative moment.

\begin{equation}
    \overline{(m^{2}\epsilon\sigma^{3})}_{xk}=2m_{k}\sum_{j}x_{j}m_{j}\left(\frac{\epsilon_{kj}}{kT}\right)\sigma_{kj}^{3}
    \label{eq:m2epssigma3_xk}
\end{equation}

% EqID: m2eps2sigma3_xk
% Source: \cite{Gross2001}, Eq.~(A.40) (appendix sequence inference)
% Status: Manual literature match
% Description: Provides a residual Helmholtz-energy relation for dispersion contribution.
% Change note: Mapped by appendix sequence as the second dispersion composition derivative moment.

\begin{equation}
    (\overline{m^2\epsilon^2\sigma^3})_{xk}=2m_{k}\sum_{j}x_{j}m_{j}\left(\frac{\epsilon_{kj}}{kT}\right)^2\sigma_{kj}^3
    \label{eq:m2eps2sigma3_xk}
\end{equation}

% EqID: c1_xk
% Source: \cite{Gross2001}, Eq.~(A.41)
% Status: Manual literature match
% Description: Provides a residual Helmholtz-energy relation for dispersion contribution.
% Change note: Mapped manually to the composition derivative of $C_1$.

\begin{equation}
    C_{1,xk}=C_{2}\zeta_{3,xk}-C_{1}^{2}\left\{m_{k}\frac{8\eta-2\eta^{2}}{\left(1-\eta\right)^{4}}-m_{k}\frac{20\eta-27\eta^{2}+12\eta^{3}-2\eta^{4}}{\left[(1-\eta)(2-\eta)\right]^{2}}\right\}
    \label{eq:c1_xk}
\end{equation}

% EqID: i1_xk
% Source: \cite{Gross2001}, Eq.~(A.42)
% Status: Close literature match
% Description: Provides a residual Helmholtz-energy relation for dispersion contribution.
% Change note: High textual similarity to a tagged equation in the cited local paper export.

\begin{equation}
    I_{1,xk}=\sum_{i=0}^{6}[a_{\mathrm{i}}(\bar{m})i\zeta_{3,xk}\eta^{i-1}+a_{\mathrm{i,x}k}\eta^{i}]
    \label{eq:i1_xk}
\end{equation}

% EqID: i2_xk
% Source: \cite{Gross2001}, Eq.~(A.43)
% Status: Close literature match
% Description: Provides a residual Helmholtz-energy relation for dispersion contribution.
% Change note: High textual similarity to a tagged equation in the cited local paper export.

\begin{equation}
    I_{2,xk}=\sum_{i=0}^{6}[b_{\mathrm{i}}(\bar{m})i\zeta_{3,xk}\eta^{i-1}+b_{\mathrm{i,xk}}\eta^{i}]
    \label{eq:i2_xk}
\end{equation}

% EqID: a_i_xk
% Source: \cite{Gross2001}, Eq.~(A.44)
% Status: Close literature match
% Description: Provides a supporting relation used in dispersion contribution.
% Change note: High textual similarity to a tagged equation in the cited local paper export.

\begin{equation}
    a_{\mathrm{i},xk}=\frac{m_{k}}{\bar{m}^{2}}a_{1i}+\frac{m_{k}}{\bar{m}^{2}}\left(3-\frac{4}{\bar{m}}\right)a_{2i}
    \label{eq:a_i_xk}
\end{equation}

% EqID: b_i_xk
% Source: \cite{Gross2001}, Eq.~(A.45)
% Status: Close literature match
% Description: Provides a supporting relation used in dispersion contribution.
% Change note: High textual similarity to a tagged equation in the cited local paper export.

\begin{equation}
    b_{\mathrm{i},xk}=\frac{m_{k}}{\bar{m}^{2}}b_{1i}+\frac{m_{k}}{\bar{m}^{2}}\left(3-\frac{4}{\bar{m}}\right)b_{2i}
    \label{eq:b_i_xk}
\end{equation}

\subsection{Association Contribution}

See~\cite{Chapman1990}.
% EqID: mu_assoc
% Source: \cite{Chapman1990} (exact equation number not available in local PDFs)
% Status: No direct numbered source in local corpus
% Description: Gives a chemical-potential contribution in association contribution.
% Change note: Association chemical-potential decomposition written in implementation form; exact numbered Chapman equation unavailable in local PDFs.

\begin{equation}
    \begin{aligned}
        \tilde{\mu}^{assoc}_{k}
        &=\tilde{a}^{\mathrm{assoc}} + Z^{assoc}
        +\left(\frac{\partial\tilde{a}^{\mathrm{assoc}}}{\partial x_{k}}\right)_{T,\nu,x_{i\neq k}}
        \\
        &\quad -\sum_{j=1}^{N}\left[x_{j}\left(\frac{\partial\tilde{a}^{\mathrm{assoc}}}{\partial x_{j}}\right)_{T,\nu,x_{i\neq j}}\right]
    \end{aligned}
    \label{eq:mu_assoc}
\end{equation}

% EqID: dares_assoc_dxk
% Source: \cite{Chapman1990} (exact equation number not available in local PDFs)
% Status: No direct numbered source in local corpus
% Description: Provides a differential relation needed for association contribution calculations.
% Change note: Composition derivative of association Helmholtz contribution; no direct numbered local Chapman source available.

\begin{equation}
    \left(\frac{\partial \tilde a^{assoc}}{\partial x_{k}}\right)_{T,v,x_{i\neq k}}
    =
    \sum_{A_{k}}\left(\ln X^{A_{k}}-\frac{X^{A_{k}}}{2}+\frac{1}{2}\right)
    +
    \sum_{i} x_{i}\sum_{A_{i}}
    \left(\frac{1}{X^{A_{i}}}-\frac{1}{2}\right)
    \left(\frac{\partial X^{A_{i}}}{\partial x_{k}}\right)_{T,v,x_{i\neq k}}
    \label{eq:dares_assoc_dxk}
\end{equation}

% EqID: dx_assoc_dxk
% Source: \cite{Chapman1990} (exact equation number not available in local PDFs)
% Status: No direct numbered source in local corpus
% Description: Provides a differential relation needed for association contribution calculations.
% Change note: Site-fraction composition derivative closure from SAFT association machinery; local progression PDFs do not provide this as a numbered equation.

\begin{equation}
    \begin{aligned}
        \left(\frac{\partial X^{A_{i}}}{\partial x_{k}}\right)_{T,v,x_{j\neq k}}
        &=
        -(X^{A_{i}})^{2}\,\rho
        \Bigg[
        \sum_{B_{k}}X^{B_{k}}\Delta^{A_{i}B_{k}}
        \\
        &\quad +
        \sum_{j}x_{j}\sum_{B_{j}}
        \left(
        \Delta^{A_{i}B_{j}}\left(\frac{\partial X^{B_{j}}}{\partial x_{k}}\right)_{T,v,x_{j\neq k}}
        +
        X^{B_{j}}\left(\frac{\partial \Delta^{A_{i}B_{j}}}{\partial x_{k}}\right)_{T,v,x_{j\neq k}}
        \right)
        \Bigg]
    \end{aligned}
    \label{eq:dx_assoc_dxk}
\end{equation}

% EqID: ddelta_assoc_dxk
% Source: \cite{Chapman1990} (exact equation number not available in local PDFs)
% Status: No direct numbered source in local corpus
% Description: Defines the Debye screening quantity used in association contribution.
% Change note: Composition derivative of association strength term; no exact numbered Chapman equation available locally.

\begin{equation}
    \left(\frac{\partial \Delta^{A_{i}B_{j}}}{\partial x_{k}}\right)_{T,v,x_{j\neq k}}
    =
    d_{ij}^{3}\kappa^{A_{i}B_{j}}
    \left[\exp\!\left(\frac{\epsilon^{A_{i}B_{j}}}{kT}\right)-1\right]
    \left(\frac{\partial g_{ij}^{hs}}{\partial x_{k}}\right)_{T,\rho,x_{j\neq k}}
    \label{eq:ddelta_assoc_dxk}
\end{equation}

\subsection{Debye and Huckel Electrolyte Term Contribution}



% EqID: mu_dh
% Source: \cite{Gross2001}, Eq.~(A.33)
% Status: Adapted implementation form
% Description: Gives a chemical-potential contribution in debye and huckel electrolyte term contribution.
% Change note: Chemical-potential contribution written via the generic Eq. (A.33)-style decomposition for the Debye-Huckel term.

\begin{equation}
    \begin{aligned}
        \tilde{\mu}^{DH}_{k}
        &=\tilde{a}^{\mathrm{DH}} + Z^{DH}
        +\left(\frac{\partial\tilde{a}^{\mathrm{DH}}}{\partial x_{k}}\right)_{T,\nu,x_{i\neq k}}
        \\
        &\quad -\sum_{j=1}^{N}\left[x_{j}\left(\frac{\partial\tilde{a}^{\mathrm{DH}}}{\partial x_{j}}\right)_{T,\nu,x_{i\neq j}}\right]
    \end{aligned}
    \label{eq:mu_dh}
\end{equation}

The original 2005 version from~\cite{Cameretti2005}

% EqID: mu_dh_2005
% Source: \cite{Cameretti2005}, Eq.~(11)
% Status: Original literature form
% Description: Defines the Debye screening quantity used in debye and huckel electrolyte term contribution.
% Change note: Baseline 2005 formulation retained for comparison to updated derivatives.

\begin{equation}
    \tilde{\mu}^{DH}_{i}= -\frac{q_{i}^2 \kappa}{24 \pi k T \epsilon}\left[2 \chi_{i}+\frac{\sum_{j} x_{j} q_{j}^2 \sigma_{k}}{\sum_{j} x_{j} q_{j}^2}\right]
    \label{eq:mu_dh_2005}
\end{equation}

% EqID: sigma_dh_2005
% Source: \cite{Cameretti2005}, Eq.~(20)
% Status: Close literature match
% Description: Defines the Debye screening quantity used in debye and huckel electrolyte term contribution.
% Change note: High textual similarity to a tagged equation in the cited local paper export.

\begin{equation}
    \sigma_{k}=\left(\frac{\partial\left(\kappa \chi_{k}\right)}{\partial \kappa}\right)_{T, \mathrm{~N}}=-2 \chi_{k}+\frac{3}{1+\kappa d_{k}}
    \label{eq:sigma_dh_2005}
\end{equation}

The updated version for concentration dependent dielectric constant from~\cite{Bulow2019}

% EqID: dares_dh_dxi
% Source: \cite{Bulow2019} (equation number unresolved in local corpus)
% Status: Literature update
% Description: Defines the Debye screening quantity used in debye and huckel electrolyte term contribution.
% Change note: Changed from the 2005 constant-dielectric form to the 2019 concentration-dependent dielectric derivative form.

\begin{equation}
    \begin{aligned}
        \left(\frac{\partial \tilde{a}^{DH}}{\partial x_{i}}\right)_{T,v,x_{j\neq i}}
        =
        -\frac{e^{2}}{12\pi\varepsilon_{0}k_{B}T}
        \Bigg[
          & \frac{1}{\varepsilon_{r}}\left(\frac{\partial\kappa}{\partial x_{i}}\right)_{T,v,x_{j\neq i}}
            \sum_{j}x_{j}z_{j}^{2}\chi_{j}
        \\
        - & \frac{\kappa}{\varepsilon_{r}^{2}}
            \left(\frac{\partial\varepsilon_{r}}{\partial x_{i}}\right)_{T,v,x_{j\neq i}}
            \sum_{j}x_{j}z_{j}^{2}\chi_{j}
        \\
        + & \frac{\kappa}{\varepsilon_{r}}
            \sum_{j}\chi_{j}z_{j}^{2}
            \left(\frac{\partial x_{j}}{\partial x_{i}}\right)_{T,v,x_{k\neq i}}
        \\
        + & \frac{\kappa}{\varepsilon_{r}}
            \sum_{j}x_{j}z_{j}^{2}
            \left(\frac{\partial\chi_{j}}{\partial x_{i}}\right)_{T,v,x_{k\neq i}}
            \Bigg]
    \end{aligned}
    \label{eq:dares_dh_dxi}
\end{equation}

% EqID: dkappa_dh_dxi
% Source: \cite{Bulow2019}, Eq.~(13)
% Status: Manual literature match
% Description: Defines the Debye screening quantity used in debye and huckel electrolyte term contribution.
% Change note: Mapped manually to the concentration-dependent dielectric derivative of Debye screening parameter.

\begin{equation}
    \begin{aligned}
         & \left(\frac{\partial\kappa}{\partial x_{i}}\right)_{T,v,x_{j\neq i}}=\frac12\left(\frac{\rho_{N}e^{2}}{k_{B}T\varepsilon_{0}\varepsilon_{r}}\sum_{j}x_{j}z_{j}^{2}\right)^{-\frac12} \\&\left\{-\frac{\rho_{N}e^2}{k_{B}T\left(\varepsilon_{0}\varepsilon_{r}\right)^2}\varepsilon_{0}\frac{\partial\left(\varepsilon_{r}\right)}{\partial x_{i}}\sum_{j}x_{j}z_{j}^2+\frac{\rho_{N}e^2}{k_{B}T\varepsilon_{0}\varepsilon_{r}}\alpha_{i}z_{i}^2\right\}\\&=\left(\frac{\rho_{N}e^{2}}{k_{B}T}\right)^{\frac{1}{2}}\left\{-\frac{1}{2}\left(\varepsilon_{0}\varepsilon_{r}\right)^{-\frac{3}{2}}\varepsilon_{0}\frac{\partial\left(\varepsilon_{r}\right)}{\partial x_{i}}\left[\sum_{j}x_{j}z_{j}^{2}\right]^{\frac{1}{2}}\right.\\&+\frac{1}{2\sqrt{\varepsilon_{0}\varepsilon_{r}}}\Bigg[\sum_{j}x_{j}z_{j}^2\Bigg]^{-\frac{1}{2}}\alpha_{i}z_{i}^2\Bigg\}
    \end{aligned}
    \label{eq:dkappa_dh_dxi}
\end{equation}

% EqID: dchi_dh_dxi
% Source: \cite{Bulow2019}, Eq.~(14)
% Status: Manual literature match
% Description: Defines the Debye screening quantity used in debye and huckel electrolyte term contribution.
% Change note: Mapped manually to the concentration derivative of the chi-function.

\begin{equation}
    \begin{aligned}
        \left(\frac{\partial\chi_{i}}{\partial x_{i}}\right)_{T,v,x_{j\neq i}} & =-\frac{9}{\left(\kappa d_{i}\right)^{4}}\biggl[\frac{3}{2}+\ln(1+\kappa d_{i})-2(1+\kappa d_{i}) \\&+\frac{1}{2}\left(1+\kappa d_{i}\right)^{2}\bigg]d_{i}\frac{\partial\kappa}{\partial x_{i}}+\frac{3}{\left(\kappa d_{i}\right)^{3}}\bigg[\frac{1}{1+\kappa d_{i}}-1+\kappa d_{i}\bigg]\\d_{i}\frac{\partial\kappa}{\partial x_{i}}&= 3\frac{\partial\kappa}{\partial x_{i}}\left\{-\frac{\chi_{i}}{\kappa}+\frac{d_{i}}{\left(\kappa d_{i}\right)^{3}}\biggl[\frac{1}{1+\kappa d_{i}}-1+\kappa d_{i}\biggr]\right\}
    \end{aligned}
    \label{eq:dchi_dh_dxi}
\end{equation}

\subsubsection{Bjerrum Treatment}

See~\cite{Bulow2021}.
% EqID: dares_dh_dxi_bjerrum
% Source: \cite{Bulow2021}, Eq.~(21)
% Status: Manual literature match
% Description: Defines the Debye screening quantity used in debye and huckel electrolyte term contribution.
% Change note: Mapped manually to the Bjerrum-treatment composition derivative of reduced Debye-Huckel Helmholtz energy.

\begin{equation}
    \begin{aligned}
        \left(\frac{\partial \tilde{a}^{DH}}{\partial x_{i}}\right)_{T,v,x_{j\neq i}}
        =
        -\frac{e^{2}}{12\pi\varepsilon_{0}k_{B}T}
        \Bigg[
          & \frac{1}{\varepsilon_{r}}\left(\frac{\partial\kappa}{\partial x_{i}}\right)_{T,v,x_{j\neq i}}
            \sum_{j}\alpha_{j}x_{j}z_{j}^{2}\chi_{j}
        \\
        - & \frac{\kappa}{\varepsilon_{r}^{2}}
            \left(\frac{\partial\varepsilon_{r}}{\partial x_{i}}\right)_{T,v,x_{j\neq i}}
            \sum_{j}\alpha_{j}x_{j}z_{j}^{2}\chi_{j}
        \\
        + & \frac{\kappa}{\varepsilon_{r}}\alpha_{i}z_{i}^{2}\chi_{i}
        \\
        + & \frac{\kappa}{\varepsilon_{r}}
            \sum_{j}\alpha_{j}x_{j}z_{j}^{2}
            \left(\frac{\partial\chi_{j}}{\partial x_{i}}\right)_{T,v,x_{k\neq i}}
            \Bigg]
    \end{aligned}
    \label{eq:dares_dh_dxi_bjerrum}
\end{equation}


% EqID: dkappa_dh_dxi_bjerrum
% Source: \cite{Bulow2021}, Eq.~(22)
% Status: Adapted notation
% Description: Defines the Debye screening quantity used in debye and huckel electrolyte term contribution.
% Change note: Moderate-to-high similarity; notation/arrangement appears adapted from the cited equation.

\begin{equation}
    \begin{aligned}
         & \left(\frac{\partial\kappa}{\partial x_{i}}\right)_{T,v,x_{j\neq i}}=\frac12\left(\frac{\rho_{N}e^{2}}{k_{B}T\varepsilon_{0}\varepsilon_{r}}\sum_{j}\alpha_{j}x_{j}z_{j}^{2}\right)^{-\frac12} \\&\left\{-\frac{\rho_{N}e^2}{k_{B}T\left(\varepsilon_{0}\varepsilon_{r}\right)^2}\varepsilon_{0}\frac{\partial\left(\varepsilon_{r}\right)}{\partial x_{i}}\sum_{j}\alpha_{j}x_{j}z_{j}^2+\frac{\rho_{N}e^2}{k_{B}T\varepsilon_{0}\varepsilon_{r}}\alpha_{i}z_{i}^2\right\}\\&=\left(\frac{\rho_{N}e^{2}}{k_{B}T}\right)^{\frac{1}{2}}\left\{-\frac{1}{2}\left(\varepsilon_{0}\varepsilon_{r}\right)^{-\frac{3}{2}}\varepsilon_{0}\frac{\partial\left(\varepsilon_{r}\right)}{\partial x_{i}}\left[\sum_{j}\alpha_{j}x_{j}z_{j}^{2}\right]^{\frac{1}{2}}\right.\\&+\frac{1}{2\sqrt{\varepsilon_{0}\varepsilon_{r}}}\Bigg[\sum_{j}\alpha_{j}x_{j}z_{j}^2\Bigg]^{-\frac{1}{2}}\alpha_{i}z_{i}^2\Bigg\}
    \end{aligned}
    \label{eq:dkappa_dh_dxi_bjerrum}
\end{equation}

% EqID: dchi_dh_dxi_bjerrum
% Source: \cite{Bulow2021}, Eq.~(23)
% Status: Manual literature match
% Description: Defines the Debye screening quantity used in debye and huckel electrolyte term contribution.
% Change note: Mapped manually to the Bjerrum-treatment chi-function derivative relation.

\begin{equation}
    \begin{aligned}
        \left(\frac{\partial\chi_{i}}{\partial x_{i}}\right)_{T,v,x_{j\neq i}} & =-\frac{9}{\left(\kappa R_{i}\right)^{4}}\biggl[\frac{3}{2}+\ln(1+\kappa R_{i})-2(1+\kappa R_{i}) \\&+\frac{1}{2}\left(1+\kappa R_{i}\right)^{2}\bigg]R_{i}\frac{\partial\kappa}{\partial x_{i}}+\frac{3}{\left(\kappa R_{i}\right)^{3}}\bigg[\frac{1}{1+\kappa R_{i}}-1+\kappa R_{i}\bigg]\\R_{i}\frac{\partial\kappa}{\partial x_{i}}&= 3\frac{\partial\kappa}{\partial x_{i}}\left\{-\frac{\chi_{i}}{\kappa}+\frac{R_{i}}{\left(\kappa R_{i}\right)^{3}}\biggl[\frac{1}{1+\kappa R_{i}}-1+\kappa R_{i}\biggr]\right\}
    \end{aligned}
    \label{eq:dchi_dh_dxi_bjerrum}
\end{equation}

\subsubsection{Dielectric Constant}

Mole-Fraction
% EqID: depsr_dxi_mole
% Source: \cite{Ascani2021}, Eq.~(11) (derived differential)
% Status: Manual literature match
% Description: Specifies dielectric-property mixing or derivative form for debye and huckel electrolyte term contribution.
% Change note: Direct derivative of the mole-fraction dielectric mixing rule.

\begin{equation}
    \left(\frac{\partial \varepsilon_{r}}{\partial x_{i}}\right)_{x_{j\neq i}}
    = \varepsilon_{r, i}
    \label{eq:depsr_dxi_mole}
\end{equation}

Mass-Fraction
% EqID: depsr_dxi_mass
% Source: \cite{Ascani2021}, Eq.~(12) (derived differential)
% Status: Manual literature match
% Description: Specifies dielectric-property mixing or derivative form for debye and huckel electrolyte term contribution.
% Change note: Direct derivative of the mass-fraction dielectric mixing rule.

\begin{equation}
    \left(\frac{\partial \varepsilon_{r}}{\partial x_{i}}\right)_{x_{j\neq i}}
    =
    \frac{MW_{i}}{\overline{MW}}\left(\varepsilon_{r,i}-\varepsilon_{r}\right)
    \label{eq:depsr_dxi_mass}
\end{equation}

Combo
% EqID: depsr_dxi_combo
% Source: \cite{Ascani2021}, Eq.~(13) (derived differential)
% Status: Manual literature match
% Description: Specifies dielectric-property mixing or derivative form for debye and huckel electrolyte term contribution.
% Change note: Direct derivative of the mixed solvent-plus-ion dielectric mixing rule.

\begin{equation}
    \left(\frac{\partial \varepsilon_{r}}{\partial x_{i}}\right)_{x_{j\neq i}}
    =
    \varepsilon_{r,sol}^{(w)}
    +
    x_{sol}\,\frac{MW_{i}}{\overline{MW}_{sol}}
    \left(\varepsilon_{r,i}-\varepsilon_{r,sol}^{(w)}\right),
    \qquad i\in\mathcal S
    \label{eq:depsr_dxi_combo}
\end{equation}

New Mixing Rule
% EqID: epsr_mix_suppressed
% Source: \cite{Bulow2021} (equation number unresolved in local corpus)
% Status: New literature extension
% Description: Provides a supporting relation used in debye and huckel electrolyte term contribution.
% Change note: Changed from earlier dielectric-mixing rules to the 2025 ion-suppression mixing form.

\begin{equation}
    \varepsilon_{r,\mathrm{mix}}(\mathbf{x})
    =
    \frac{\varepsilon_{sf}(\mathbf{x})}{1+7.01\,x_{\mathrm{ion}}(\mathbf{x})},
    \qquad
    \varepsilon_{sf}\equiv \varepsilon_{r,\mathrm{solvent,mix}}^{\mathrm{salt\text{-}free}} .
    \label{eq:epsr_mix_suppressed}
\end{equation}

% EqID: depsr_mix_dxi
% Source: \cite{Figiel2025}, Eq.~(11)-Eq.~(12) (derived differential)
% Status: Manual literature match
% Description: Provides a differential relation needed for debye and huckel electrolyte term contribution calculations.
% Change note: Derivative of the 2025 dielectric mixing rule via chain rule.

\begin{equation}
    \left(\frac{\partial \varepsilon_{r,\mathrm{mix}}}{\partial x_{i}}\right)_{T,v,x_{j\neq i}}
    =
    \frac{1}{1+7.01\,x_{\mathrm{ion}}}
    \left(\frac{\partial \varepsilon_{sf}}{\partial x_{i}}\right)_{T,v,x_{j\neq i}}
    -\frac{7.01\,\varepsilon_{sf}}{\left(1+7.01\,x_{\mathrm{ion}}\right)^2}
    \left(\frac{\partial x_{\mathrm{ion}}}{\partial x_{i}}\right)_{T,v,x_{j\neq i}} .
    \label{eq:depsr_mix_dxi}
\end{equation}

% EqID: epsr_sf
% Source: \cite{Figiel2025}, Eq.~(12)
% Status: Manual literature match
% Description: Provides a supporting relation used in debye and huckel electrolyte term contribution.
% Change note: Mapped manually to the salt-free solvent dielectric mixture definition.

\begin{equation}
    \varepsilon_{sf}(\mathbf{x})
    =
    \sum_{s\in\mathcal{S}} w_{s}^{sf}\,\varepsilon_{r,s},
    \qquad
    w_{s}^{sf}
    =
    \frac{x_{s} M_{s}}{\displaystyle \sum_{m\in\mathcal{S}} x_{m} M_{m}},
    \qquad
    D\equiv \sum_{m\in\mathcal{S}} x_{m} M_{m} .
    \label{eq:epsr_sf}
\end{equation}

% EqID: depsr_sf_dxi
% Source: \cite{Figiel2025}, Eq.~(12) (derived differential)
% Status: Manual literature match
% Description: Provides a differential relation needed for debye and huckel electrolyte term contribution calculations.
% Change note: Derivative of solvent-mixture dielectric expression in Eq. (12).

\begin{equation}
    \left(\frac{\partial \varepsilon_{sf}}{\partial x_{i}}\right)_{T,v,x_{j\neq i}}
    =
    \sum_{s\in\mathcal{S}} \varepsilon_{r,s}
    \left(\frac{\partial w_{s}^{sf}}{\partial x_{i}}\right)_{T,v,x_{j\neq i}}
    =
    \begin{cases}
        \dfrac{M_{i}}{D}\left(\varepsilon_{r,i}-\varepsilon_{sf}\right), & i\in\mathcal{S},    \\[10pt]
        0,                                                             & i\notin\mathcal{S}.
    \end{cases}
    \label{eq:depsr_sf_dxi}
\end{equation}

% EqID: x_ion_total
% Source: \cite{Figiel2025}, Eq.~(11) (derived helper)
% Status: Manual literature match
% Description: Provides a differential relation needed for debye and huckel electrolyte term contribution calculations.
% Change note: Ion-fraction helper used in the Eq. (11) differential form.

\begin{equation}
    x_{\mathrm{ion}}(\mathbf{x})=\sum_{m\in\mathcal{I}} x_{m},
    \qquad
    \left(\frac{\partial x_{\mathrm{ion}}}{\partial x_{i}}\right)_{T,v,x_{j\neq i}}
    =
    \begin{cases}
        1, & i\in\mathcal{I},    \\
        0, & i\notin\mathcal{I}.
    \end{cases}
    \label{eq:x_ion_total}
\end{equation}

% EqID: depsr_mix_dxi_piecewise
% Source: \cite{Figiel2025}, Eq.~(11)-Eq.~(12) (derived closed form)
% Status: Manual literature match
% Description: Provides a differential relation needed for debye and huckel electrolyte term contribution calculations.
% Change note: Piecewise closed-form derivative obtained by combining Eq. (11) and Eq. (12).

\begin{equation}
    \left(\frac{\partial \varepsilon_{r,\mathrm{mix}}}{\partial x_{i}}\right)_{T,v,x_{j\neq i}}
    =
    \begin{cases}
        \dfrac{1}{1+7.01\,x_{\mathrm{ion}}}\;
        \dfrac{M_{i}}{\displaystyle \sum_{m\in\mathcal{S}} x_{m} M_{m}}
        \left(\varepsilon_{r,i}-\varepsilon_{sf}\right),
         & i\in\mathcal{S}, \\[14pt]
        -\dfrac{7.01\,\varepsilon_{sf}}{\left(1+7.01\,x_{\mathrm{ion}}\right)^{2}},
         & i\in\mathcal{I}.
    \end{cases}
    \label{eq:depsr_mix_dxi_piecewise}
\end{equation}



\subsection{Born Electrolyte Term Contribution}

% EqID: mu_born
% Source: \cite{Gross2001}, Eq.~(A.33) and \cite{Bulow2021a}, Eq.~(3)
% Status: Manual literature match
% Description: Gives a chemical-potential contribution in born electrolyte term contribution.
% Change note: Born chemical-potential contribution written as Eq. (A.33)-style decomposition using the Part I Born composition derivative term.

\begin{equation}
    \begin{aligned}
        \tilde{\mu}^{Born}_{k}
        &=\tilde{a}^{\mathrm{Born}} + Z^{Born}
        +\left(\frac{\partial\tilde{a}^{\mathrm{Born}}}{\partial x_{k}}\right)_{T,\nu,x_{i\neq k}}
        \\
        &\quad -\sum_{j=1}^{N}\left[x_{j}\left(\frac{\partial\tilde{a}^{\mathrm{Born}}}{\partial x_{j}}\right)_{T,\nu,x_{i\neq j}}\right]
    \end{aligned}
    \label{eq:mu_born}
\end{equation}


Version 1:
How the reference ``ePC-SAFT advanced, Part~I\@: Physical meaning of including a concentration-dependent dielectric constant in the born term and in the Debye-H\"uckel theory'' has it (equations 3)~\cite{Bulow2021a}
% EqID: dares_born_dxi_v1
% Source: \cite{Bulow2021a}, Eq.~(3)
% Status: Adapted notation
% Description: Specifies dielectric-property mixing or derivative form for born electrolyte term contribution.
% Change note: Moderate-to-high similarity; notation/arrangement appears adapted from the cited equation.

\begin{equation}
    \begin{aligned}\left(\frac{\partial \tilde{a}^{Born}}{\partial x_{i}}\right)_{T,v_{N},x_{j\neq i}}=-\frac{e^{2}}{4\pi k_{B}T\varepsilon_{0}}\biggl[\left(1-\frac{1}{\varepsilon_{r}}\right)\frac{z_{i}^{2}}{d^{\text{Born}}_{i}} + \left(\frac{1}{\varepsilon_{r}^{2}}\right) \left(\frac{\partial\varepsilon_{r}}{\partial x_{i}}\right)_{T,v,x_{j\neq i}} \biggr]\end{aligned}
    \label{eq:dares_born_dxi_v1}
\end{equation}

Version 2:
Potentially Fixed Differential from the paper
% EqID: dares_born_dxi_v2
% Source: \cite{Bulow2021a}, Eq.~(3)
% Status: Project-specific modification
% Description: Specifies dielectric-property mixing or derivative form for born electrolyte term contribution.
% Change note: Changed from Version 1 differential by restoring the composition summation in the dielectric-derivative term.

\begin{equation}
    \begin{aligned}\left(\frac{\partial \tilde{a}^{Born}}{\partial x_{i}}\right)_{T,v_{N},x_{j\neq i}}=-\frac{e^{2}}{4\pi k_{B}T\varepsilon_{0}}\biggl[\left(1-\frac{1}{\varepsilon_{r}}\right)\frac{z_{i}^{2}}{d^{\text{Born}}_{i}} + \left(\frac{1}{\varepsilon_{r}^{2}}\right)\sum_{j} \frac{x_{j} z_{j}^2}{d^{\text{Born}}_{j}} \left(\frac{\partial\varepsilon_{r}}{\partial x_{i}}\right)_{T,v,x_{j\neq i}} \biggr]\end{aligned}
    \label{eq:dares_born_dxi_v2}
\end{equation}

\subsubsection{Solvation Shell Model + Dielectric Saturation}

See~\cite{Figiel2025}.
Version 3

% EqID: dares_born_ssmds_dxi_v3
% Source: \cite{Figiel2025}, Eq.~(6)
% Status: Project-specific modification
% Description: Specifies dielectric-property mixing or derivative form for born electrolyte term contribution.
% Change note: Changed from earlier Born differential variants toward the SSM+DS implementation path.

\begin{equation}
    \begin{aligned}
        \left(\frac{\partial \tilde a_{\mathrm{SSM+DS}}^{\mathrm{Born}}}{\partial x_{k}}\right)_{T,v,x_{j\neq k}}
        =
        -\frac{e^{2}}{4\pi\varepsilon_{0}k_{\mathrm{B}}T}
        \Bigg[
         & \left(1-\frac{1}{\varepsilon_{r}}\right)\frac{z_{k}^{2}}{d_{k}^{\mathrm{Born}}+\Delta d_{k}}
        \\
         & +\left(1-\frac{1}{\varepsilon_{r,\mathrm{ion}}}\right) z_{k}^{2}
            \left(\frac{1}{d_{k}^{\mathrm{Born}}}-\frac{1}{d_{k}^{\mathrm{Born}}+\Delta d_{k}}\right)
        \\
         & +\left(\frac{1}{\varepsilon_{r}^{2}}\right)
            \left(\frac{\partial \varepsilon_{r}}{\partial x_{k}}\right)_{T,v,x_{j\neq k}}
            \left(\sum_{i} x_{i} z_{i}^{2}\frac{1}{d_{i}^{\mathrm{Born}}+\Delta d_{i}}\right)
            \Bigg].
    \end{aligned}
    \label{eq:dares_born_ssmds_dxi_v3}
\end{equation}

Version 4
% EqID: dares_born_dxi_v4
% Source: \cite{Figiel2025} (equation number unresolved in local corpus)
% Status: Project-specific modification
% Description: Provides a differential relation needed for born electrolyte term contribution calculations.
% Change note: Changed from Version 3 to the final refactored derivative form with explicit $D_i$ and $\Delta d_m$ terms.

\begin{equation}
    \begin{aligned}
        \left(\frac{\partial a^{\mathrm{Born}}}{\partial x_{i}}\right)_{T,v,x_{j\neq i}}
         & =
        -\frac{e^2}{4\pi\varepsilon_{0} k_{B} T}
        \Bigg[
            z_{i}^2\!\left(
            \left(1-\frac{1}{\varepsilon_{r,\mathrm{ion}}}\right)\frac{1}{D_{i}}
            +
            \left(1-\frac{1}{\varepsilon_{r,\mathrm{solv}}}\right)\left(\frac{1}{d_{i}^{\mathrm{Born}}}-\frac{1}{D_{i}}\right)
            \right)
        \\
         & \qquad\qquad
            +\sum_{j} x_{j} z_{j}^2
            \Bigg(
            \left(\frac{1}{\varepsilon_{r,\mathrm{ion}}}-\frac{1}{\varepsilon_{r,\mathrm{solv}}}\right)
            \frac{1}{D_{j}^2}
            \frac{\partial \Delta d_{m}}{\partial x_{i}}
            \\
         & \qquad\qquad\qquad\qquad
            +
            \frac{1}{\varepsilon_{r,\mathrm{solv}}^2}
            \frac{\partial \varepsilon_{r,\mathrm{solv}}}{\partial x_{i}}
            \left(\frac{1}{d_{j}^{\mathrm{Born}}}-\frac{1}{D_{j}}\right)
            \Bigg)
            \Bigg],
    \end{aligned}
    \label{eq:dares_born_dxi_v4}
\end{equation}

% EqID: ddelta_d_dxi
% Source: \cite{Figiel2025}, Eq.~(6)
% Status: Adapted implementation form
% Description: Provides a differential relation needed for born electrolyte term contribution calculations.
% Change note: Lower similarity; likely algebraically adapted for implementation or combined terms.

\begin{equation}
    \frac{\partial \Delta d_{m}}{\partial x_{i}}=\frac{d_{m}^{\text {Born }}}{\left|z_{m}\right|} \frac{\partial f_{\text {mix }}}{\partial x_{i}}=\frac{d_{m}^{\text {Born }}}{\left|z_{m}\right|} f_{i},
    \label{eq:ddelta_d_dxi}
\end{equation}

% EqID: d_born_effective
% Source: \cite{Figiel2025}, Eq.~(8)
% Status: Adapted implementation form
% Description: Provides a supporting relation used in born electrolyte term contribution.
% Change note: Lower similarity; likely algebraically adapted for implementation or combined terms.

\begin{equation}
    D_{m}=d_{m}^{\text {Born }}+\Delta d_{m}, \Delta d_{m}=\frac{\left(f_{\text {mix }}-1\right)}{\left|z_{m}\right|} d_{m}^{\text {Born }} .
    \label{eq:d_born_effective}
\end{equation}

\section{Fugacity Coefficient}

See~\cite{Gross2001}.
% EqID: lnphi_total
% Source: \cite{Gross2001}, Eq.~(A.32)
% Status: Manual literature match
% Description: Gives the total fugacity-coefficient relation in fugacity coefficient.
% Change note: Mapped manually to the residual-chemical-potential form used in the PC-SAFT appendix.

\begin{equation}
    \ln\varphi_{k}
    =
    \tilde{\mu}_{k}^{\mathrm{res}}-\ln Z
    =
    \frac{\mu_{k}^{\mathrm{res}}}{kT}-\ln Z
    =
    \sum_\alpha \ln \varphi^\alpha
    =
    \sum_\alpha \mu^\alpha+\sum_\alpha\left[-\frac{Z^\alpha}{Z-1} \ln Z\right]
    \label{eq:lnphi_total}
\end{equation}

% EqID: lnphi_total_sum
% Source: Project decomposition based on Eq.~\eqref{eq:lnphi_total}
% Status: Project-specific organization
% Description: Gives the explicit fugacity-coefficient decomposition in fugacity coefficient.
% Change note: Written in explicit non-summation form to match the contribution-by-contribution structure used throughout this section.

\begin{equation}
    \ln\varphi_{k}
    =
    \ln\varphi_{k}^{hc}
    +
    \ln\varphi_{k}^{disp}
    +
    \ln\varphi_{k}^{assoc}
    +
    \ln\varphi_{k}^{DH}
    +
    \ln\varphi_{k}^{Born}
    \label{eq:lnphi_total_sum}
\end{equation}

% EqID: lnphi_alpha
% Source: Project decomposition based on Eq.~\eqref{eq:lnphi_total}
% Status: Project-specific organization
% Description: Gives the generic contribution fugacity-coefficient relation in fugacity coefficient.
% Change note: Uses only the explicit $Z^\alpha$ allocation requested for the contribution-resolved fugacity-coefficient terms.

\begin{equation}
    \ln\varphi_{k}^{\alpha}
    =
    \tilde{\mu}_{k}^{\alpha}
    -
    \frac{Z^{\alpha}}{Z-1}\ln Z
    \label{eq:lnphi_alpha}
\end{equation}

% EqID: lnphi_alpha_near_ideal
% Source: Derived from Eq.~\eqref{eq:lnphi_alpha}
% Status: Derived approximation
% Description: Gives the near-ideal approximation for a contribution fugacity coefficient in fugacity coefficient.
% Change note: Retained explicitly as an approximation only, using the $Z\rightarrow 1$ limit requested for documentation.

\begin{equation}
    \ln\varphi_{k}^{\alpha}
    \approx
    \tilde{\mu}_{k}^{\alpha}
    -
    Z^{\alpha}
    \label{eq:lnphi_alpha_near_ideal}
\end{equation}

\subsection{Hard-Chain Reference Contribution}

% EqID: lnphi_hc
% Source: Project decomposition based on Eq.~\eqref{eq:lnphi_alpha}
% Status: Project-specific organization
% Description: Gives the hard-chain fugacity-coefficient contribution in fugacity coefficient.
% Change note: Hard-chain specialization of the explicit $Z^\alpha$ fugacity-coefficient split.

\begin{equation}
    \ln\varphi_{k}^{hc}
    =
    \tilde{\mu}_{k}^{hc}
    -
    \frac{Z^{hc}}{Z-1}\ln Z
    \label{eq:lnphi_hc}
\end{equation}

\subsection{Dispersion Contribution}

% EqID: lnphi_disp
% Source: Project decomposition based on Eq.~\eqref{eq:lnphi_alpha}
% Status: Project-specific organization
% Description: Gives the dispersion fugacity-coefficient contribution in fugacity coefficient.
% Change note: Dispersion specialization of the explicit $Z^\alpha$ fugacity-coefficient split.

\begin{equation}
    \ln\varphi_{k}^{disp}
    =
    \tilde{\mu}_{k}^{disp}
    -
    \frac{Z^{disp}}{Z-1}\ln Z
    \label{eq:lnphi_disp}
\end{equation}

\subsection{Association Contribution}

% EqID: lnphi_assoc
% Source: Project decomposition based on Eq.~\eqref{eq:lnphi_alpha}
% Status: Project-specific organization
% Description: Gives the association fugacity-coefficient contribution in fugacity coefficient.
% Change note: Association specialization of the explicit $Z^\alpha$ fugacity-coefficient split.

\begin{equation}
    \ln\varphi_{k}^{assoc}
    =
    \tilde{\mu}_{k}^{assoc}
    -
    \frac{Z^{assoc}}{Z-1}\ln Z
    \label{eq:lnphi_assoc}
\end{equation}

\subsection{Debye and Huckel Electrolyte Term Contribution}

% EqID: lnphi_dh
% Source: Project decomposition based on Eq.~\eqref{eq:lnphi_alpha}
% Status: Project-specific organization
% Description: Gives the Debye-Huckel fugacity-coefficient contribution in fugacity coefficient.
% Change note: Debye-Huckel specialization of the explicit $Z^\alpha$ fugacity-coefficient split.

\begin{equation}
    \ln\varphi_{k}^{DH}
    =
    \tilde{\mu}_{k}^{DH}
    -
    \frac{Z^{DH}}{Z-1}\ln Z
    \label{eq:lnphi_dh}
\end{equation}

\subsection{Born Electrolyte Term Contribution}

% EqID: lnphi_born
% Source: Project decomposition based on Eq.~\eqref{eq:lnphi_alpha}
% Status: Project-specific organization
% Description: Gives the Born fugacity-coefficient contribution in fugacity coefficient.
% Change note: Born specialization of the explicit $Z^\alpha$ fugacity-coefficient split.

\begin{equation}
    \ln\varphi_{k}^{Born}
    =
    \tilde{\mu}_{k}^{Born}
    -
    \frac{Z^{Born}}{Z-1}\ln Z
    \label{eq:lnphi_born}
\end{equation}

\section{Activity Coefficient}

See~\cite{Gross2001,Figiel2025}.
% EqID: gamma_sym
% Source: Standard thermodynamic definition
% Status: Project-specific organization
% Description: Gives the symmetric-reference activity-coefficient definition in activity coefficient.
% Change note: Added explicitly in non-log form so the activity section mirrors the fugacity-coefficient organization.

\begin{equation}
    \gamma_{i}
    =
    \frac{\varphi_{i}(T,p,\mathbf{x})}{\varphi_{0i}(T,p,x_{i}=1)}
    \label{eq:gamma_sym}
\end{equation}

% EqID: gamma_asym_inf
% Source: Standard thermodynamic definition
% Status: Project-specific organization
% Description: Gives the infinite-dilution activity-coefficient definition in activity coefficient.
% Change note: Added explicitly in non-log form for ions and solutes referenced to infinite dilution.

\begin{equation}
    \gamma_{i}^{*}
    =
    \frac{\varphi_{i}(T,p,\mathbf{x})}{\varphi_{i}^{\infty}(T,p,x_{i}\to 0)}
    \label{eq:gamma_asym_inf}
\end{equation}

% EqID: lngamma_sym
% Source: Derived from Eq.~\eqref{eq:gamma_sym}
% Status: Derived relation
% Description: Gives the symmetric-reference logarithmic activity-coefficient definition in activity coefficient.
% Change note: Expressed only in terms of fugacity coefficients, as requested for the activity section.

\begin{equation}
    \ln\gamma_{i}
    =
    \ln\varphi_{i}
    -
    \ln\varphi_{0i}
    \label{eq:lngamma_sym}
\end{equation}

% EqID: lngamma_asym_inf
% Source: Derived from Eq.~\eqref{eq:gamma_asym_inf}
% Status: Derived relation
% Description: Gives the infinite-dilution logarithmic activity-coefficient definition in activity coefficient.
% Change note: Expressed only in terms of fugacity coefficients, as requested for the activity section.

\begin{equation}
    \ln\gamma_{i}^{*}
    =
    \ln\varphi_{i}
    -
    \ln\varphi_{i}^{\infty}
    \label{eq:lngamma_asym_inf}
\end{equation}

% EqID: lngamma_asym_sum
% Source: Project decomposition based on Eq.~\eqref{eq:lngamma_asym_inf}
% Status: Project-specific organization
% Description: Gives the explicit infinite-dilution activity-coefficient decomposition in activity coefficient.
% Change note: Written in explicit non-summation form to mirror the fugacity-coefficient decomposition.

\begin{equation}
    \ln\gamma_{k}^{*}
    =
    \ln\gamma_{k}^{hc,*}
    +
    \ln\gamma_{k}^{disp,*}
    +
    \ln\gamma_{k}^{assoc,*}
    +
    \ln\gamma_{k}^{DH,*}
    +
    \ln\gamma_{k}^{Born,*}
    \label{eq:lngamma_asym_sum}
\end{equation}

% EqID: gamma_alpha_asym
% Source: Project decomposition based on Eq.~\eqref{eq:gamma_asym_inf}
% Status: Project-specific organization
% Description: Gives the generic contribution activity-coefficient definition in activity coefficient.
% Change note: Defined only from contribution fugacity coefficients, not from chemical-potential expressions.

\begin{equation}
    \gamma_{k}^{\alpha,*}
    =
    \frac{\varphi_{k}^{\alpha}}{\varphi_{k}^{\alpha,\infty}}
    \label{eq:gamma_alpha_asym}
\end{equation}

% EqID: lngamma_alpha_asym
% Source: Derived from Eq.~\eqref{eq:gamma_alpha_asym}
% Status: Derived relation
% Description: Gives the generic logarithmic contribution activity-coefficient definition in activity coefficient.
% Change note: Written only in terms of contribution fugacity coefficients, consistent with the requested section logic.

\begin{equation}
    \ln\gamma_{k}^{\alpha,*}
    =
    \ln\varphi_{k}^{\alpha}
    -
    \ln\varphi_{k}^{\alpha,\infty}
    \label{eq:lngamma_alpha_asym}
\end{equation}

% EqID: gamma_pm_asym
% Source: Standard thermodynamic definition
% Status: Project-specific organization
% Description: Gives the mean-ionic infinite-dilution activity coefficient in activity coefficient.
% Change note: Added in non-log form to parallel the logarithmic mean-ionic relation already used in electrolyte work.

\begin{equation}
    \gamma_{\pm}^{*}
    =
    \left(
    \left(\gamma_{c}^{*}\right)^{\nu_{c}}
    \left(\gamma_{a}^{*}\right)^{\nu_{a}}
    \right)^{\frac{1}{\nu_{c}+\nu_{a}}}
    \label{eq:gamma_pm_asym}
\end{equation}

% EqID: lngamma_pm_asym
% Source: Derived from Eq.~\eqref{eq:gamma_pm_asym}
% Status: Derived relation
% Description: Gives the logarithmic mean-ionic infinite-dilution activity coefficient in activity coefficient.
% Change note: Written explicitly in the standard stoichiometric-weighted logarithmic form.

\begin{equation}
    \ln\gamma_{\pm}^{*}
    =
    \frac{\nu_{c}\ln\gamma_{c}^{*}+\nu_{a}\ln\gamma_{a}^{*}}{\nu_{c}+\nu_{a}}
    \label{eq:lngamma_pm_asym}
\end{equation}

% EqID: lngamma_pm_alpha_asym
% Source: Project decomposition based on Eq.~\eqref{eq:lngamma_pm_asym}
% Status: Project-specific organization
% Description: Gives the logarithmic contribution mean-ionic activity coefficient in activity coefficient.
% Change note: Contribution form written directly from the contribution activity-coefficient terms.

\begin{equation}
    \ln\gamma_{\pm}^{\alpha,*}
    =
    \frac{\nu_{c}\ln\gamma_{c}^{\alpha,*}+\nu_{a}\ln\gamma_{a}^{\alpha,*}}{\nu_{c}+\nu_{a}}
    \label{eq:lngamma_pm_alpha_asym}
\end{equation}

\subsection{Hard-Chain Reference Contribution}

% EqID: lngamma_hc_asym
% Source: Project decomposition based on Eq.~\eqref{eq:lngamma_alpha_asym}
% Status: Project-specific organization
% Description: Gives the hard-chain activity-coefficient contribution in activity coefficient.
% Change note: Hard-chain specialization of the infinite-dilution contribution activity-coefficient definition.

\begin{equation}
    \ln\gamma_{k}^{hc,*}
    =
    \ln\varphi_{k}^{hc}
    -
    \ln\varphi_{k}^{hc,\infty}
    \label{eq:lngamma_hc_asym}
\end{equation}

\subsection{Dispersion Contribution}

% EqID: lngamma_disp_asym
% Source: Project decomposition based on Eq.~\eqref{eq:lngamma_alpha_asym}
% Status: Project-specific organization
% Description: Gives the dispersion activity-coefficient contribution in activity coefficient.
% Change note: Dispersion specialization of the infinite-dilution contribution activity-coefficient definition.

\begin{equation}
    \ln\gamma_{k}^{disp,*}
    =
    \ln\varphi_{k}^{disp}
    -
    \ln\varphi_{k}^{disp,\infty}
    \label{eq:lngamma_disp_asym}
\end{equation}

\subsection{Association Contribution}

% EqID: lngamma_assoc_asym
% Source: Project decomposition based on Eq.~\eqref{eq:lngamma_alpha_asym}
% Status: Project-specific organization
% Description: Gives the association activity-coefficient contribution in activity coefficient.
% Change note: Association specialization of the infinite-dilution contribution activity-coefficient definition.

\begin{equation}
    \ln\gamma_{k}^{assoc,*}
    =
    \ln\varphi_{k}^{assoc}
    -
    \ln\varphi_{k}^{assoc,\infty}
    \label{eq:lngamma_assoc_asym}
\end{equation}

\subsection{Debye and Huckel Electrolyte Term Contribution}

% EqID: lngamma_dh_asym
% Source: Project decomposition based on Eq.~\eqref{eq:lngamma_alpha_asym}
% Status: Project-specific organization
% Description: Gives the Debye-Huckel activity-coefficient contribution in activity coefficient.
% Change note: Debye-Huckel specialization of the infinite-dilution contribution activity-coefficient definition.

\begin{equation}
    \ln\gamma_{k}^{DH,*}
    =
    \ln\varphi_{k}^{DH}
    -
    \ln\varphi_{k}^{DH,\infty}
    \label{eq:lngamma_dh_asym}
\end{equation}

\subsection{Born Electrolyte Term Contribution}

% EqID: lngamma_born_asym
% Source: Project decomposition based on Eq.~\eqref{eq:lngamma_alpha_asym}
% Status: Project-specific organization
% Description: Gives the Born activity-coefficient contribution in activity coefficient.
% Change note: Born specialization of the infinite-dilution contribution activity-coefficient definition.

\begin{equation}
    \ln\gamma_{k}^{Born,*}
    =
    \ln\varphi_{k}^{Born}
    -
    \ln\varphi_{k}^{Born,\infty}
    \label{eq:lngamma_born_asym}
\end{equation}


\section{Enthalpy, Entropy, and Gibbs Free Energy}

\subsection{Entropy}

% EqID: h_res
% Source: \cite{Figiel2025}, Eq.~(13)
% Status: Adapted implementation form
% Description: Provides a differential relation needed for entropy calculations.
% Change note: Lower similarity; likely algebraically adapted for implementation or combined terms.

\begin{equation}
    \tilde{h}^{\mathrm{res}} = \frac{\hat{h}^{\mathrm{res}}}{RT}=-T\left(\frac{\partial\tilde{a}^{\mathrm{res}}}{\partial T}\right)_{\rho,x_{i}}+(Z-1)
    \label{eq:h_res}
\end{equation}

\subsection{Entropy}

% EqID: s_res_from_s_vol
% Source: \cite{Gross2001}, Eq.~(A.47)
% Status: Manual literature match
% Description: Provides a supporting relation used in entropy.
% Change note: Mapped manually to the residual-entropy relation with logarithmic compressibility correction.

\begin{equation}
    \tilde{s}^{\mathrm{res}} = \frac{\hat{s}^{\mathrm{res}}(P,T)}{R}=\frac{\hat{s}^{\mathrm{res}}(\nu,T)}{R}+\ln(Z)
    \label{eq:s_res_from_s_vol}
\end{equation}

% EqID: s_res
% Source: \cite{Gross2001}, Eq.~(A.48)
% Status: Manual literature match
% Description: Provides a differential relation needed for entropy calculations.
% Change note: Mapped manually to the residual-entropy temperature-derivative form.

\begin{equation}
    \tilde{s}^{\mathrm{res}} = \frac{\hat{s}^{\mathrm{res}}(P,T)}{R}=-T\left[\left(\frac{\partial\tilde{a}^{\mathrm{res}}}{\partial T}\right)_{\rho,x_{i}}+\frac{\tilde{a}^{\mathrm{res}}}{T}\right]+\ln(Z)
    \label{eq:s_res}
\end{equation}

\subsection{Gibbs Free Energy}

% EqID: g_res_from_hs
% Source: \cite{Gross2001}, Eq.~(A.49)
% Status: Manual literature match
% Description: Provides a supporting relation used in gibbs free energy.
% Change note: Mapped manually to the residual Gibbs relation via enthalpy and entropy.

\begin{equation}
    \tilde{g}^{\mathrm{res}}=\frac{\hat{g}^{\mathrm{res}}}{RT}=\frac{\hat{h}^{\mathrm{res}}}{RT}-\frac{\hat{s}^{\mathrm{res}}(P,T)}{R}
    \label{eq:g_res_from_hs}
\end{equation}

% EqID: g_res_from_ares
% Source: \cite{Gross2001}, Eq.~(A.50)
% Status: Manual literature match
% Description: Provides a residual Helmholtz-energy relation for gibbs free energy.
% Change note: Mapped manually to the residual Gibbs relation in Helmholtz/compressibility form.

\begin{equation}
    \tilde{g}^{\mathrm{res}}=\tilde{a}^{\mathrm{res}}+(Z-1)-\ln(Z)
    \label{eq:g_res_from_ares}
\end{equation}

Gibbs Energy of Solvation

% EqID: delta_g_solv_inf_x
% Source: \cite{Figiel2025}, Eq.~(14)
% Status: Adapted implementation form
% Description: Provides a residual Helmholtz-energy relation for gibbs free energy.
% Change note: Lower similarity; likely algebraically adapted for implementation or combined terms.

\begin{equation}
    \Delta^{\mathrm{solv},x}G_{i}^{\infty}(T,p,x_{j\neq i},x_{i}\to0)=RT\ln(\varphi_{i}^{\infty}(T,p,x_{j\neq i},x_{i}\to0))
    \label{eq:delta_g_solv_inf_x}
\end{equation}

Gibbs Energy of Transfer at Infinite Dilution from solvent S1 to solvent S2
% EqID: delta_g_transfer_inf
% Source: \cite{Figiel2025}, Eq.~(14)
% Status: Adapted implementation form
% Description: Provides a residual Helmholtz-energy relation for gibbs free energy.
% Change note: Lower similarity; likely algebraically adapted for implementation or combined terms.

\begin{equation}
    \Delta^{\mathrm{tr}}G_{i}^{\infty,\mathrm{S}1\to\mathrm{S}2}(T,p,x_{j\neq i},x_{i}\rightarrow0)=RT\ln\left(\frac{\varphi_{i}^{\infty,\mathrm{S}2}}{\varphi_{i}^{\infty,\mathrm{S}1}}\right)
    \label{eq:delta_g_transfer_inf}
\end{equation}

\subsection{Temperature Differential Equations}

% EqID: dares_dT
% Source: \cite{Gross2001}, Eq.~(A.51)
% Status: Project-specific extension
% Description: Provides a differential relation needed for temperature differential equations calculations.
% Change note: Mapped to Eq. (A.51) and extended here by including association, Debye-Huckel, and Born derivatives.

\begin{equation}
    \left(\frac{\partial\tilde{a}^\mathrm{res}}{\partial T}\right)_{\rho,x_{i}} =\left(\frac{\partial\tilde{a}^\mathrm{hc}}{\partial T}\right)_{\rho,x_{i}} +\left(\frac{\partial\tilde{a}^\mathrm{disp}}{\partial T}\right)_{\rho,x_{i}}
    +\left(\frac{\partial\tilde{a}^\mathrm{assoc}}{\partial T}\right)_{\rho,x_{i}}
    +\left(\frac{\partial\tilde{a}^\mathrm{DH}}{\partial T}\right)_{\rho,x_{i}}
    +\left(\frac{\partial\tilde{a}^\mathrm{Born}}{\partial T}\right)_{\rho,x_{i}}
    \label{eq:dares_dT}
\end{equation}

\subsubsection{Hard-Chain Reference Contribution}

% EqID: dares_hc_dT
% Source: \cite{Gross2001}, Eq.~(A.54)
% Status: Manual literature match
% Description: Provides a differential relation needed for temperature differential equations calculations.
% Change note: Mapped manually to the hard-chain temperature derivative expression.

\begin{equation}
    \left(\frac{\partial\tilde{a}^{\mathrm{hc}}}{\partial T}\right)_{\rho,x_{i}}=\bar{m}\left(\frac{\partial\tilde{a}^{\mathrm{hs}}}{\partial T}\right)_{\rho,x_{i}}-\sum_{i}x_{\mathrm{i}}(m_{\mathrm{i}}-1)(g_{ii}^{\mathrm{hs}})^{-1}\left(\frac{\partial g_{ii}^{\mathrm{hs}}}{\partial T}\right)_{\rho,x_{i}}
    \label{eq:dares_hc_dT}
\end{equation}

% EqID: dares_hs_dT
% Source: \cite{Gross2001}, Eq.~(A.55)
% Status: Manual literature match
% Description: Provides a differential relation needed for temperature differential equations calculations.
% Change note: Mapped manually to the hard-sphere temperature derivative expression.

\begin{equation}
    \begin{aligned}
        \left(\frac{\partial\tilde{a}^{hs}}{\partial T}\right)_{\rho,x_{i}}
        &=\frac{1}{\zeta_{0}}\Bigg[
        \frac{3(\zeta_{1,T}\zeta_{2}+\zeta_{1}\zeta_{2,T})}{(1-\zeta_{3})}
        +\frac{3\zeta_{1}\zeta_{2}\zeta_{3,T}}{(1-\zeta_{3})^{2}}
        \\
        &\quad
        +\frac{3{\zeta_{2}}^{2}{\zeta_{2,T}}}{\zeta_{3}{(1-\zeta_{3})}^{2}}
        +\frac{{\zeta_{2}}^{3}{\zeta_{3,T}}(3{\zeta_{3}}-1)}{{\zeta_{3}}^{2}{(1-\zeta_{3})}^{3}}
        \\
        &\quad
        +\left(\frac{3{\zeta_{2}}^{2}{\zeta_{2,T}}{\zeta_{3}}-2{\zeta_{2}}^{3}{\zeta_{3,T}}}{{\zeta_{3}}^{3}}\right)\ln(1-{\zeta_{3}})
        +\left({\zeta_{0}}-\frac{{\zeta_{2}}^{3}}{{\zeta_{3}}^{2}}\right)\frac{\zeta_{3,T}}{(1-\zeta_{3})}
        \Bigg]
    \end{aligned}
    \label{eq:dares_hs_dT}
\end{equation}

% EqID: dg_hs_dT
% Source: \cite{Gross2001}, Eq.~(A.57)
% Status: Manual literature match
% Description: Provides a differential relation needed for temperature differential equations calculations.
% Change note: Mapped manually to the contact-value temperature derivative expression.

\begin{equation}
    \begin{aligned}
        \frac{\partial g_{ii}^{\mathrm{hs}}}{\partial T}
        &=\frac{\zeta_{3,T}}{\left(1-\zeta_{3}\right)^{2}}
        +\left(\frac{1}{2}d_{i,T}\right)\frac{3\zeta_{2}}{\left(1-\zeta_{3}\right)^{2}}
        \\
        &\quad +\left(\frac{1}{2}d_{i}\right)\left(\frac{3\zeta_{2,T}}{\left(1-\zeta_{3}\right)^{2}}+\frac{6\zeta_{2}\zeta_{3,T}}{\left(1-\zeta_{3}\right)^{3}}\right)
        +\left(\frac{1}{2}d_{i}d_{i,T}\right)\frac{2\zeta_{2}^{2}}{\left(1-\zeta_{3}\right)^{3}}
        \\
        &\quad +\left(\frac{1}{2}d_{i}\right)^{2}\left(\frac{4\xi_{2}\xi_{2,T}}{\left(1-\xi_{3}\right)^{3}}+\frac{6\xi_{2}{}^{2}\xi_{3,T}}{\left(1-\xi_{3}\right)^{4}}\right)
    \end{aligned}
    \label{eq:dg_hs_dT}
\end{equation}

% EqID: d_segment_dT
% Source: \cite{Gross2001}, Eq.~(A.9) (derived temperature differential)
% Status: Manual literature match
% Description: Provides a differential relation needed for temperature differential equations calculations.
% Change note: Computed by differentiating Eq. (A.9) with respect to temperature.

\begin{equation}
    d_{i,T}=\frac{\partial d_{i}}{\partial T}=\sigma_{i}\left(3\frac{\epsilon_{i}}{kT^2}\right)\left[-0.12\exp\left(-3\frac{\epsilon_{i}}{kT}\right)\right]
    \label{eq:d_segment_dT}
\end{equation}

% EqID: zeta_n_dT
% Source: \cite{Gross2001}, Eq.~(A.53)
% Status: Manual literature match
% Description: Provides a differential relation needed for temperature differential equations calculations.
% Change note: Mapped manually to the temperature derivative of zeta moments.

\begin{equation}
    \zeta_{n,T}=\frac{\partial\zeta_{n}}{\partial T}=\frac{\pi}{6}\rho\sum_{i}x_{i}m_{i}nd_{i,T}\left(d_{i}\right)^{n-1}\quad n\in\{1,2,3\}
    \label{eq:zeta_n_dT}
\end{equation}

% EqID: half_d_identity
% Source: \cite{Figiel2025}, Eq.~(20)
% Status: Adapted implementation form
% Description: Provides a supporting relation used in temperature differential equations.
% Change note: Lower similarity; likely algebraically adapted for implementation or combined terms.

\begin{equation}
    \frac{1}{2}d_{i}=\left(\frac{d_{i}d_{i}}{d_{i}+d_{i}}\right)
    \label{eq:half_d_identity}
\end{equation}

\subsubsection{Dispersion Contribution}

% EqID: dares_disp_dT
% Source: \cite{Gross2001}, Eq.~(A.58)
% Status: Manual literature match
% Description: Provides a differential relation needed for temperature differential equations calculations.
% Change note: Mapped manually to the dispersion temperature derivative expression.

\begin{equation}
    \begin{aligned}
        \left(\frac{\partial\tilde{a}^{\mathrm{disp}}}{\partial T}\right)_{\rho,x_{i}} & =-2\pi\rho\left(\frac{\partial I_{1}}{\partial T}-\frac{I_{1}}{T}\right)\overline{m^{2}\epsilon\sigma^{3}}- \\&\pi\rho\bar{m}\left[\frac{\partial C_{1}}{\partial T}I_{2}+C_{1}\frac{\partial I_{2}}{\partial T}-2C_{1}\frac{I_{2}}{T}\right]\overline{m^{2}\epsilon^{2}\sigma^{3}}
    \end{aligned}
    \label{eq:dares_disp_dT}
\end{equation}

\subsubsection{Debye and Huckel Electrolyte Term Contribution}

% EqID: dares_dh_dT
% Source: \cite{Cameretti2005}, Eq.~(13)-Eq.~(14) (derived temperature differential)
% Status: Derived from literature equation
% Description: Defines the Debye screening quantity used in temperature differential equations.
% Change note: Temperature differential form derived from Debye-Huckel term and chi-function definitions.

\begin{equation}
    \frac{d \tilde{a}^{D H}}{d T}=-\frac{1}{12 \pi k_{B} \varepsilon_{0} \varepsilon_{r}} \sum_{i} x_{i}\left(z_{i} e\right)^2\left[\left(-2 \chi_{i}+\frac{3}{1+\kappa a_{i}}\right)\left(-\frac{\kappa}{2 T^2}\right)-\frac{\kappa \chi_{i}}{T^2}\right]
    \label{eq:dares_dh_dT}
\end{equation}

\subsubsection{Born Contribution}

% EqID: dares_born_dT
% Source: \cite{Bulow2021a}, Eq.~(2) (derived temperature differential)
% Status: Derived from literature equation
% Description: Specifies dielectric-property mixing or derivative form for temperature differential equations.
% Change note: Temperature derivative of the Born contribution derived from the Part I Born Helmholtz equation.

\begin{equation}
    \left(\frac{\partial\tilde{a}^{\mathrm{Born}}}{\partial T}\right)_{\rho,x_{i}}=\frac{e^2}{4 \pi \varepsilon_{0} k_{B} T^2}\left(1-\frac{1}{\varepsilon_{r}}\right) \sum_{i} \frac{x_{i} z_{i}^2}{a_{i}}
    \label{eq:dares_born_dT}
\end{equation}

\section{Salt Basis Conversions}

% EqID: gamma_pm_molality
% Source: \cite{Figiel2025}, Eq.~(22)
% Status: Close literature match
% Description: Gives an activity or fugacity-coefficient relation in salt basis conversions.
% Change note: High textual similarity to a tagged equation in the cited local paper export.

\begin{equation}
    \gamma_{\pm}^{*,m}=\frac{\gamma_{\pm}^{*,x}}{1+M_{\mathrm{solvent}}\cdot\tilde{m}_{\mathrm{solute}}\cdot\sum_{i}\nu_{i,\mathrm{ion}}}
    \label{eq:gamma_pm_molality}
\end{equation}



\printbibliography\end{document}
